\documentclass[11pt]{article}
\usepackage[margin=1in]{geometry}
\usepackage[T1]{fontenc}
\usepackage[utf8]{inputenc}
\usepackage[french]{babel}
\usepackage{amsmath,amssymb,amsthm}
\usepackage[colorlinks=true,linkcolor=blue,urlcolor=blue]{hyperref}
\newtheorem{problem}{Problème}
\newenvironment{refs}{\par\small\noindent\emph{Références :}\begin{itemize}\setlength\itemsep{0pt}}{\end{itemize}\normalsize}
\title{Hors de la Caverne \\ \Large Géométrie}
\author{}
\date{}
\begin{document}
\maketitle
\noindent\emph{Des problèmes, pas de réponses.} Chaque problème ci-dessous est posé
sans solution. Les références indiquent où trouver les connaissances nécessaires.
\medskip


\section*{Collège}

\begin{problem}[{Pourquoi les angles d'un triangle font 180$^\circ$ --- Collège}]
\textbf{GEO-41.} \textbf{Problème.} Soit $ ABC $ un triangle. Tracez la droite passant par $ A $ et parallèle à
$ (BC) $.
\begin{enumerate}
\item[(a)] En utilisant les angles alternes-internes formés par cette parallèle avec $ (AB) $ puis avec
$ (AC) $, démontrez que
\[ \widehat{A} + \widehat{B} + \widehat{C} = 180^\circ . \]
Nommez explicitement, à chaque étape, la propriété du cours que vous employez.
\item[(b)] Déduisez-en que la somme des angles d'un quadrilatère non croisé vaut $ 360^\circ $, puis que
celle d'un polygone convexe à $ n $ côtés vaut $ (n-2) \times 180^\circ $. Justifiez le découpage
que vous utilisez.
\item[(c)] Sur une sphère, on peut tracer un triangle dont les trois angles sont droits : partez du pôle
Nord, descendez jusqu'à l'équateur, parcourez un quart d'équateur, remontez. Sa somme d'angles vaut
$ 270^\circ $. Expliquez précisément quelle étape de votre démonstration de (a) devient fausse dans
ce cas.
\end{enumerate}

\end{problem}

\noindent C'est la proposition 32 du livre I des \textit{Éléments} d'Euclide, et la question (c)
en désigne le point sensible : la démonstration ne tient que parce qu'on a pu tracer \textit{une seule}
parallèle à $ (BC) $ passant par $ A $. Cet énoncé est le cinquième postulat d'Euclide, qu'il
avait lui-même isolé comme suspect et que vingt siècles de mathématiciens ont tenté en vain de
déduire des quatre autres. Gauss, Bolyai et Lobatchevski ont fini par comprendre, dans les années
1820-1830, qu'on ne le pouvait pas : il existe des géométries cohérentes où il est faux, et où la
somme des angles d'un triangle n'est jamais $ 180^\circ $. Autrement dit, ce résultat de sixième est
équivalent au postulat le plus discuté de l'histoire des mathématiques --- le problème GEO-06, plus bas,
reprend l'affaire.

\begin{refs}
\item Contexte : Somme des angles d'un triangle --- Wikipédia (\url{https://fr.wikipedia.org/wiki/Somme_des_angles_d%27un_triangle})
\item Contexte : Angles alternes-internes --- Wikipédia (\url{https://fr.wikipedia.org/wiki/Angles_alternes-internes})
\item Contexte : Géométrie sphérique --- Wikipédia (\url{https://fr.wikipedia.org/wiki/G%C3%A9om%C3%A9trie_sph%C3%A9rique})
\end{refs}

\bigskip

\begin{problem}[{Pythagore par découpage --- Collège}]
\textbf{GEO-42.} \textbf{Problème.} Soit un triangle rectangle de côtés de l'angle droit $ a $ et $ b $, et
d'hypoténuse $ c $. Prenez deux carrés identiques de côté $ a + b $. Dans le premier, placez
quatre copies du triangle dans les quatre coins de façon à laisser au centre un quadrilatère de côté
$ c $ ; dans le second, placez les quatre mêmes copies de façon à laisser deux carrés vides, de
côtés $ a $ et $ b $.
\begin{enumerate}
\item[(a)] Démontrez que le quadrilatère central de la première figure est bien un carré. C'est le point que
tout le monde saute : il faut prouver que ses angles sont droits, en utilisant que les angles aigus du
triangle rectangle ont pour somme $ 90^\circ $.
\item[(b)] En comparant les aires des deux grands carrés, démontrez que
\[ a^2 + b^2 = c^2 . \]
\item[(c)] Le triangle de côtés $ 3, 4, 5 $ est rectangle. Cela découle-t-il du théorème que vous venez de
démontrer, ou d'un autre énoncé ? Formulez précisément l'énoncé nécessaire et dites en quoi il diffère
du précédent.
\item[(d)] Un camarade affirme : \og{} j'ai vérifié sur vingt triangles rectangles dessinés et mesurés, donc le
théorème est vrai. \fg{} Expliquez pourquoi ce n'est pas une démonstration, et ce que ses mesures
établissent réellement.
\end{enumerate}

\end{problem}

\noindent Le résultat était connu et utilisé bien avant Pythagore : la tablette babylonienne
Plimpton 322, vers 1800 av. J.-C., aligne des triplets d'entiers vérifiant la relation, et les cordes
à nœuds des arpenteurs égyptiens s'en servaient pour tracer des angles droits. La démonstration par
découpage demandée ici est essentiellement celle qui accompagne le \textit{Zhoubi Suanjing} chinois, où
la figure porte le nom de \textit{xian tu}, et on la retrouve chez Bhāskara II au XIIe siècle
en Inde, accompagnée, dit la tradition, du seul mot \og{} Regarde ! \fg{}. C'est le théorème qui compte le
plus de démonstrations différentes --- plusieurs centaines recensées, dont une due à un futur président
des États-Unis, James Garfield, en 1876. La question (c) sépare deux énoncés qu'on confond presque
toujours : le théorème et sa réciproque, et seule la réciproque permet de \textit{conclure} qu'un
triangle est rectangle.

\begin{refs}
\item Contexte : Théorème de Pythagore --- Wikipédia (\url{https://fr.wikipedia.org/wiki/Th%C3%A9or%C3%A8me_de_Pythagore})
\item Contexte : Plimpton 322 --- Wikipédia (\url{https://fr.wikipedia.org/wiki/Plimpton_322})
\item Contexte : Zhoubi Suanjing --- Wikipédia (\url{https://fr.wikipedia.org/wiki/Zhoubi_Suanjing})
\item Contexte : Preuve sans mots --- Wikipédia (\url{https://fr.wikipedia.org/wiki/Preuve_sans_mots})
\end{refs}

\bigskip

\begin{problem}[{La médiatrice, ensemble des points équidistants --- Collège}]
\textbf{GEO-43.} \textbf{Problème.} Soient $ A $ et $ B $ deux points distincts, $ I $ le milieu de
$ [AB] $, et $ (d) $ la perpendiculaire à $ (AB) $ passant par $ I $.
\begin{enumerate}
\item[(a)] Démontrez que tout point $ M $ de $ (d) $ vérifie $ MA = MB $.
\item[(b)] Démontrez la réciproque : tout point $ M $ du plan vérifiant $ MA = MB $ appartient à
$ (d) $. C'est cette moitié-là qui justifie qu'on parle de \textit{l'}ensemble des points
équidistants et non seulement d'une famille de points équidistants.
\item[(c)] Soient $ A $, $ B $, $ C $ trois points non alignés. En appliquant (a) et (b) aux
médiatrices de $ [AB] $ et de $ [BC] $, démontrez que les trois médiatrices du triangle
$ ABC $ se coupent en un même point, et qu'il existe un cercle passant par les trois sommets.
\item[(d)] Trois villages veulent creuser un puits à égale distance des trois. Démontrez qu'un tel puits
existe toujours, sauf dans un cas que vous préciserez, et expliquez géométriquement ce qui se passe
dans ce cas.
\end{enumerate}

\end{problem}

\noindent La double implication de (a) et (b) est le premier exemple, dans une scolarité, d'un
\textit{ensemble caractérisé par une propriété} --- un lieu géométrique. Le mot est ancien : les Grecs
appelaient \textit{topos} une figure définie par une condition, et la plus grande partie de la géométrie
classique consiste à reconnaître qu'une condition apparemment compliquée décrit en réalité une droite
ou un cercle. Euclide construit le cercle circonscrit à la proposition 5 du livre IV. La question (d)
est là pour l'exception : les trois médiatrices d'un triangle aplati sont parallèles, et le puits
part à l'infini --- un cas dégénéré qu'une bonne rédaction n'a pas le droit de passer sous silence.

\begin{refs}
\item Contexte : Médiatrice --- Wikipédia (\url{https://fr.wikipedia.org/wiki/M%C3%A9diatrice})
\item Contexte : Cercle circonscrit --- Wikipédia (\url{https://fr.wikipedia.org/wiki/Cercle_circonscrit})
\end{refs}

\bigskip

\begin{problem}[{Encadrer $\pi$, puis l'aire du disque --- Collège}]
\textbf{GEO-44.} \textbf{Problème.} On considère un cercle de rayon $ 1 $, donc de périmètre $ 2\pi $.
\begin{enumerate}
\item[(a)] Inscrivez dans ce cercle un carré, et circonscrivez-lui un autre carré. Calculez les deux
périmètres et démontrez l'encadrement $ 2\sqrt{2} \le \pi \le 4 $, en énonçant clairement
l'hypothèse géométrique que vous admettez sur les longueurs.
\item[(b)] Remplacez le carré inscrit par un hexagone régulier. Démontrez que son côté vaut exactement
$ 1 $ --- décomposez l'hexagone en six triangles et montrez qu'ils sont équilatéraux --- puis
déduisez-en que $ \pi \ge 3 $.
\item[(c)] Découpez maintenant le disque en un grand nombre de secteurs égaux très fins, et rangez-les tête
bêche pour former une figure proche d'un rectangle. Expliquez pourquoi les dimensions de ce
\og{} rectangle \fg{} sont proches de $ \pi $ et de $ 1 $, et pourquoi cela rend plausible que l'aire du
disque vaille $ \pi $.
\item[(d)] Dites précisément ce que l'argument de (c) ne démontre pas. Quel mot, absent du programme du
collège, faudrait-il définir pour qu'il devienne une démonstration ?
\end{enumerate}

\end{problem}

\noindent Archimède écrit vers 250 av. J.-C. un petit traité, \textit{De la mesure du cercle},
qui fait exactement cela : en partant de l'hexagone et en doublant quatre fois le nombre de côtés
jusqu'au polygone à 96 côtés, inscrit puis circonscrit, il obtient $ 223/71 < \pi < 22/7 $.
C'est le premier encadrement de l'histoire assorti d'une démonstration, et le procédé --- la méthode
d'exhaustion héritée d'Eudoxe --- est l'ancêtre direct du calcul intégral. Archimède démontre aussi, dans
le même traité, que l'aire du disque est celle d'un triangle de base le périmètre et de hauteur le
rayon, ce qui est précisément l'énoncé que le découpage de (c) fait deviner. La question (d) doit vous
mener au mot manquant : \textit{limite}. Vous le rencontrerez au lycée, et la page
Analyse réelle et théorie de la mesure (\url{pure-real-analysis.html}) montre ce qu'il a fallu de travail pour le
définir proprement --- deux mille ans après Archimède.

\begin{refs}
\item Contexte : Pi --- Wikipédia (\url{https://fr.wikipedia.org/wiki/Pi})
\item Contexte : Méthode d'exhaustion --- Wikipédia (\url{https://fr.wikipedia.org/wiki/M%C3%A9thode_d%27exhaustion})
\item Contexte : Archimède --- Wikipédia (\url{https://fr.wikipedia.org/wiki/Archim%C3%A8de})
\item Voir aussi : Analyse réelle et théorie de la mesure (\url{pure-real-analysis.html})
\end{refs}

\bigskip

\begin{problem}[{Thalès, l'ombre et la pyramide --- Collège}]
\textbf{GEO-45.} \textbf{Problème.} Un piquet vertical de $ 1{,}50 $ m projette une ombre de $ 2{,}40 $ m. Au même
instant, l'ombre d'un arbre voisin mesure $ 18 $ m.
\begin{enumerate}
\item[(a)] Faites une figure faisant apparaître la configuration de Thalès, et calculez la hauteur de
l'arbre. Justifiez chaque étape par le théorème, énoncé en entier.
\item[(b)] Cette méthode repose sur une hypothèse physique que la figure ne montre pas. Laquelle ? Expliquez
pourquoi la conclusion serait fausse si les rayons du Soleil n'étaient pas sensiblement parallèles, et
si le piquet et l'arbre n'étaient pas tous deux verticaux.
\item[(c)] La tradition rapporte que Thalès mesura ainsi la hauteur de la grande pyramide de Khéops. Le
problème est en réalité plus difficile que celui de l'arbre : le pied de la pyramide n'est pas
accessible. Expliquez quelle longueur supplémentaire il faut mesurer et pourquoi.
\item[(d)] Construisez une figure où deux triangles partagent un sommet et où les longueurs sont dans le même
rapport, mais où la conclusion du théorème est fausse parce que les droites ne sont pas parallèles.
Dites ce que cela prouve sur l'usage du théorème.
\end{enumerate}

\end{problem}

\noindent Thalès de Milet, au VIe siècle av. J.-C., est le premier personnage à qui
la tradition grecque attribue des démonstrations plutôt que des recettes --- le passage, décisif, de
\og{} cela marche \fg{} à \og{} voici pourquoi \fg{}. L'histoire de la pyramide nous vient de Diogène Laërce et de
Plutarque, huit siècles plus tard, et deux versions circulent : l'une où Thalès attend l'heure à
laquelle son ombre égale sa taille, l'autre où il utilise le rapport général, celle du présent
problème. On ne peut pas trancher, et il faut le dire ainsi. Ce qui n'est pas légendaire, c'est le
contenu du théorème : il exprime que la géométrie euclidienne est invariante par changement d'échelle,
propriété que la géométrie sphérique de GEO-41 ne possède pas --- sur une sphère, agrandir une figure la
déforme.

\begin{refs}
\item Contexte : Théorème de Thalès --- Wikipédia (\url{https://fr.wikipedia.org/wiki/Th%C3%A9or%C3%A8me_de_Thal%C3%A8s})
\item Contexte : Thalès --- Wikipédia (\url{https://fr.wikipedia.org/wiki/Thal%C3%A8s})
\item Contexte : Pyramide de Khéops --- Wikipédia (\url{https://fr.wikipedia.org/wiki/Pyramide_de_Kh%C3%A9ops})
\end{refs}

\bigskip

\begin{problem}[{Le patron d'un cube et les faces d'un dé --- Collège}]
\textbf{GEO-46.} \textbf{Problème.} Un patron de cube est une figure faite de six carrés reliés par des côtés, qui se
plie en un cube sans recouvrement.
\begin{enumerate}
\item[(a)] Trouvez le plus grand nombre possible de patrons de cube, deux patrons superposables par rotation
ou retournement étant considérés comme identiques. Il y en a exactement onze : décrivez la méthode qui
vous garantit de n'en oublier aucun.
\item[(b)] Démontrez que deux carrés d'un patron qui partagent un côté ne deviennent jamais des faces
opposées du cube.
\item[(c)] Sur un dé usuel, deux faces opposées portent des nombres dont la somme vaut $ 7 $. Prenez un
patron en croix, numérotez-en une face $ 1 $ et une face voisine $ 2 $ : montrez qu'il reste
exactement deux façons de compléter le dé, et que ces deux dés ne sont pas superposables.
\item[(d)] On double l'arête d'un cube. Démontrez que son aire est multipliée par $ 4 $ et son volume par
$ 8 $, et expliquez pourquoi les exposants $ 2 $ et $ 3 $ apparaissent ici.
\end{enumerate}

\end{problem}

\noindent Le dénombrement des onze patrons est l'un des rares problèmes de collège où l'on
doit démontrer qu'une liste est \textit{complète} --- non pas trouver des exemples, mais fermer la porte.
La bonne méthode consiste à classer les patrons par la longueur de leur plus longue bande de carrés
alignés, ce qui transforme une recherche tâtonnante en un raisonnement fini. Les deux dés de la
question (c) sont symétriques l'un de l'autre sans être superposables : les dés occidentaux et
certains dés asiatiques anciens n'ont pas la même chiralité, et un joueur attentif peut le voir. La
question (d) contient en germe toute la théorie des dimensions --- pourquoi une fourmi peut tomber sans
dommage et pas un éléphant --- et c'est le même exposant $ 3 $ qui explique que Kepler, en empilant
des boulets de canon, ait posé un problème resté ouvert jusqu'en 1998.

\begin{refs}
\item Contexte : Patron (géométrie) --- Wikipédia (\url{https://fr.wikipedia.org/wiki/Patron_(g%C3%A9om%C3%A9trie)})
\item Contexte : Cube --- Wikipédia (\url{https://fr.wikipedia.org/wiki/Cube})
\item Contexte : Dé --- Wikipédia (\url{https://fr.wikipedia.org/wiki/D%C3%A9})
\end{refs}

\bigskip

\begin{problem}[{Le plus court chemin jusqu'à la rivière --- Collège}]
\textbf{GEO-47.} \textbf{Problème.} Une droite $ (d) $ représente la berge d'une rivière. Deux points $ A $ et
$ B $ sont situés du même côté de $ (d) $. Un promeneur part de $ A $, doit toucher la berge en
un point $ M $, puis rejoindre $ B $ ; il cherche à rendre $ AM + MB $ le plus petit possible.
\begin{enumerate}
\item[(a)] Soit $ A' $ le symétrique de $ A $ par rapport à $ (d) $. Démontrez que pour tout point
$ M $ de $ (d) $, on a $ AM = A'M $, donc $ AM + MB = A'M + MB $.
\item[(b)] Démontrez, à l'aide de l'inégalité triangulaire, que $ A'M + MB $ est minimal lorsque $ M $
est le point d'intersection de $ [A'B] $ avec $ (d) $, et que ce point existe et est unique.
\item[(c)] Démontrez qu'en ce point l'angle entre $ [AM] $ et $ (d) $ est égal à l'angle entre
$ [MB] $ et $ (d) $.
\item[(d)] Reprenez le problème lorsque $ A $ et $ B $ sont de part et d'autre de $ (d) $. Que devient
la solution, et pourquoi le détour par le symétrique n'est-il alors plus nécessaire ?
\end{enumerate}

\end{problem}

\noindent Héron d'Alexandrie, au Ier siècle, pose et résout ce problème dans sa
\textit{Catoptrique}, et en tire une conclusion audacieuse : puisque la lumière réfléchie suit un chemin
dont l'angle d'incidence égale l'angle de réflexion, c'est que \textit{la lumière prend le chemin le plus
court}. C'est le premier principe de minimalité de la physique. Fermat le reprend en 1662, corrige
\og{} le plus court \fg{} en \og{} le plus rapide \fg{} et en déduit la loi de la réfraction ; de là descend toute la
mécanique de Maupertuis, d'Euler et de Lagrange, et l'idée que les lois de la nature s'écrivent comme
des minimisations. Tout cela tient dans la symétrie axiale de la question (a) : un pliage de feuille
qui remplace un problème d'optimisation par une ligne droite.

\begin{refs}
\item Contexte : Symétrie axiale --- Wikipédia (\url{https://fr.wikipedia.org/wiki/Sym%C3%A9trie_axiale})
\item Contexte : Inégalité triangulaire --- Wikipédia (\url{https://fr.wikipedia.org/wiki/In%C3%A9galit%C3%A9_triangulaire})
\item Contexte : Héron d'Alexandrie --- Wikipédia (\url{https://fr.wikipedia.org/wiki/H%C3%A9ron_d%27Alexandrie})
\item Contexte : Principe de Fermat --- Wikipédia (\url{https://fr.wikipedia.org/wiki/Principe_de_Fermat})
\end{refs}

\bigskip

\begin{problem}[{Des triangles de même aire --- Collège, short}]
\textbf{GEO-48.} \textbf{Problème.} Soient $ A $ et $ B $ deux points d'une droite $ (d) $, et
$ (d') $ une droite parallèle à $ (d) $. Démontrez que tous les triangles $ ABM $, lorsque
$ M $ parcourt $ (d') $, ont la même aire, et expliquez pourquoi cette aire ne dépend pas de la
position de $ M $ alors que la forme du triangle, elle, change du tout au tout.
\end{problem}

\noindent C'est le lemme du \og{} cisaillement \fg{} : on peut faire glisser le sommet d'un triangle
le long d'une parallèle à la base sans rien changer à l'aire, et cette seule remarque démontre en trois
lignes la moitié des formules d'aire du collège. Euclide s'en sert aux propositions 37 et 38 du livre
I, juste avant d'attaquer Pythagore, et la démonstration célèbre de la proposition 47 n'est presque
rien d'autre que deux cisaillements bien choisis.

\begin{refs}
\item Contexte : Aire d'un triangle --- Wikipédia (\url{https://fr.wikipedia.org/wiki/Aire_d%27un_triangle})
\end{refs}

\bigskip

\begin{problem}[{Quels polygones pavent le plan --- Collège}]
\textbf{GEO-49.} \textbf{Problème.} On veut recouvrir le plan, sans trou ni chevauchement, avec des copies d'un même
polygone régulier, les carreaux se joignant côté contre côté et sommet contre sommet.
\begin{enumerate}
\item[(a)] En utilisant le résultat de GEO-41, démontrez que l'angle intérieur d'un polygone régulier à
$ n $ côtés vaut
\[ \frac{(n-2) \times 180^\circ}{n} . \]
\item[(b)] Autour de chaque sommet du pavage, les angles des carreaux doivent totaliser exactement
$ 360^\circ $. Démontrez que l'angle intérieur doit donc diviser $ 360^\circ $, et déduisez-en que
seuls le triangle équilatéral, le carré et l'hexagone régulier conviennent.
\item[(c)] Traitez explicitement le cas du pentagone régulier : calculez son angle et montrez où le calcul
échoue. Dessinez trois pentagones autour d'un point et mesurez le trou qui reste.
\item[(d)] Parmi les trois pavages possibles, l'hexagone est celui qui, à aire donnée, utilise la plus petite
longueur de cloisons. Cherchez ce que les abeilles en font, et énoncez en une phrase le théorème du
nid d'abeilles en précisant qui l'a démontré et quand.
\end{enumerate}

\end{problem}

\noindent La restriction aux trois pavages réguliers est déjà chez Pappus d'Alexandrie
(IVe siècle), qui remarque au passage que les abeilles ont \og{} choisi \fg{} l'hexagone et y voit
un signe de sagesse géométrique. L'intuition qu'il exprime --- l'hexagone minimise la cloison --- est
restée une conjecture pendant seize siècles, jusqu'à ce que Thomas Hales la démontre en 1999 : c'est
le théorème du nid d'abeilles. Levez ensuite la contrainte de régularité et le sujet explose : les
pavages du plan par des figures quelconques sont classés en dix-sept groupes de symétrie, dont les
mosaïques de l'Alhambra de Grenade illustrent une bonne part six siècles avant qu'on ne sache les
compter, et la question de savoir quels pentagones convexes pavent le plan n'a été close qu'en 2017,
par Michaël Rao.

\begin{refs}
\item Contexte : Pavage du plan --- Wikipédia (\url{https://fr.wikipedia.org/wiki/Pavage_du_plan})
\item Contexte : Polygone régulier --- Wikipédia (\url{https://fr.wikipedia.org/wiki/Polygone_r%C3%A9gulier})
\item Contexte : Pavage hexagonal --- Wikipédia (\url{https://fr.wikipedia.org/wiki/Pavage_hexagonal})
\item Contexte : Théorème du nid d'abeilles --- Wikipédia (en anglais) (\url{https://en.wikipedia.org/wiki/Honeycomb_theorem})
\end{refs}

\bigskip

\begin{problem}[{Deux miroirs parallèles --- Collège, short}]
\textbf{GEO-50.} \textbf{Problème.} Soient $ (d_1) $ et $ (d_2) $ deux droites parallèles
distinctes. Démontrez que la composée des deux symétries axiales d'axes $ (d_1) $ puis $ (d_2) $
est une translation, et déterminez sa direction, son sens et sa longueur en fonction de la distance
entre les deux droites.
\end{problem}

\noindent C'est le premier théorème de composition de transformations qu'on puisse démontrer
sans machinerie, et il ouvre une porte : les symétries axiales engendrent toutes les isométries du
plan, si bien que translations et rotations ne sont pas des objets nouveaux mais des produits de
réflexions. Remplacez \og{} parallèles \fg{} par \og{} sécantes \fg{} et vous obtiendrez une rotation d'angle double
--- la même démonstration, une figure différente --- et vous serez à la porte de la théorie des groupes,
qui classe les frises et les pavages de GEO-49.

\begin{refs}
\item Contexte : Symétrie axiale --- Wikipédia (\url{https://fr.wikipedia.org/wiki/Sym%C3%A9trie_axiale})
\item Contexte : Groupe de frise --- Wikipédia (\url{https://fr.wikipedia.org/wiki/Groupe_de_frise})
\end{refs}

\bigskip

\begin{problem}[{Voir n'est pas démontrer --- Collège, short}]
\textbf{GEO-51.} \textbf{Problème.} Tracez très soigneusement, sur cinq triangles quelconques et
franchement différents, le point de concours des médianes $ G $, celui des médiatrices $ O $ et
celui des hauteurs $ H $ ; constatez à la règle que ces trois points paraissent toujours alignés.
Rédigez ensuite, en une demi-page, pourquoi ces cinq figures --- et cent de plus --- ne démontrent rien,
et ce qu'il faudrait produire à la place.
\end{problem}

\noindent L'alignement est réel : c'est la droite d'Euler, publiée en 1765, deux mille ans
après Euclide --- rappel que la géométrie du triangle n'était pas close. Sa démonstration est au-dessus
du collège et vous attend en GEO-03, plus bas sur cette page. Ce qui est à votre portée aujourd'hui,
et qui compte davantage, c'est de savoir nommer l'écart entre une régularité observée et un théorème :
la conjecture de Pólya, vérifiée sur des centaines de millions de cas, s'est révélée fausse ; celle de
Mertens aussi. Une figure vous dit où chercher, elle ne vous dit jamais que vous avez trouvé.

\begin{refs}
\item Contexte : Droite d'Euler --- Wikipédia (\url{https://fr.wikipedia.org/wiki/Droite_d%27Euler})
\item Voir aussi : L'art de la démonstration (\url{proofs.html})
\end{refs}

\bigskip

\begin{problem}[{Reconnaître deux droites parallèles --- Collège, short}]
\textbf{GEO-52.} \textbf{Problème.} Deux droites $ (d) $ et $ (d') $ sont coupées par une
sécante $ (s) $, et deux angles alternes-internes ainsi formés sont égaux. Démontrez que $ (d) $
et $ (d') $ sont parallèles, en nommant précisément l'axiome ou le théorème du cours qui autorise
chacun de vos pas.
\end{problem}

\noindent C'est la proposition 27 du livre I d'Euclide, et sa place dans les \textit{Éléments}
est remarquable : elle est démontrée \textit{sans} le cinquième postulat, alors que sa réciproque --- si
les droites sont parallèles alors les angles alternes-internes sont égaux, proposition 29 --- en a
besoin. Les deux énoncés se ressemblent comme deux gouttes d'eau et n'ont pas du tout le même statut
logique ; savoir dans quel sens on raisonne est ici la seule chose qui compte.

\begin{refs}
\item Contexte : Angles alternes-internes --- Wikipédia (\url{https://fr.wikipedia.org/wiki/Angles_alternes-internes})
\item Contexte : Axiome des parallèles --- Wikipédia (\url{https://fr.wikipedia.org/wiki/Axiome_des_parall%C3%A8les})
\end{refs}

\bigskip


\section*{Lycée}

\begin{problem}[{Le théorème de l'angle inscrit --- Lycée}]
\textbf{GEO-01.} \textbf{Problème.} Soient $ A $ et $ B $ deux points d'un cercle de centre $ O $,
et soit $ P $ un point quelconque du grand arc $ AB $.
\begin{enumerate}
\item[(a)] Démontrez que l'angle au centre vaut le double de l'angle inscrit :
\[ \angle AOB = 2\,\angle APB. \]
\item[(b)] Déduisez-en le théorème de l'angle inscrit dans un demi-cercle : un angle inscrit dans un demi-cercle est droit.
\item[(c)] Déduisez-en que les angles opposés d'un quadrilatère inscriptible ont pour somme $ 180^\circ $.
\end{enumerate}

\end{problem}

\noindent Le théorème de l'angle inscrit est la proposition 20 du livre III des
\textit{Éléments} d'Euclide, et la tradition attribue à Thalès de Milet (VIe siècle
av. J.-C.) le cas du demi-cercle --- peut-être le tout premier théorème de l'histoire attribué
à une personne nommée. C'est le moteur de la \og{} chasse aux angles \fg{} : presque tout problème
sérieux de géométrie du cercle en olympiade tourne grâce à lui. La démonstration exige une
disjonction de cas ; trouver la bonne décomposition, c'est tout l'enjeu.

\begin{refs}
\item Contexte : Angle inscrit --- Wikipédia (\url{https://fr.wikipedia.org/wiki/Angle_inscrit})
\item Contexte : Angle inscrit dans un demi-cercle --- Wikipédia (\url{https://fr.wikipedia.org/wiki/Angle_inscrit_dans_un_demi-cercle})
\item Lecture : H. S. M. Coxeter, \textit{Introduction to Geometry}, ch. 1
\end{refs}

\bigskip

\begin{problem}[{Le théorème de Ptolémée --- Lycée}]
\textbf{GEO-02.} \textbf{Problème.} Soit $ ABCD $ un quadrilatère inscrit dans un cercle.
\begin{enumerate}
\item[(a)] Démontrez le théorème de Ptolémée :
\[ AC \cdot BD = AB \cdot CD + AD \cdot BC. \]
\item[(b)] Déduisez-en le théorème de Pythagore comme cas particulier.
\item[(c)] Pour aller plus loin, démontrez la forme en inégalité : pour \textit{quatre points quelconques}
$ A, B, C, D $ du plan, $ AC \cdot BD \le AB \cdot CD + AD \cdot BC $, avec
égalité exactement lorsque $ ABCD $ est inscriptible dans cet ordre.
\end{enumerate}

\end{problem}

\noindent Claude Ptolémée a démontré cette identité dans l'\textit{Almageste}
(IIe siècle apr. J.-C.) : elle était le cœur calculatoire de sa table des cordes ---
en pratique une table des sinus construite pour l'astronomie. Appliquée à des quadrilatères
bien choisis, elle fournit les formules d'addition des sinus, et c'est ainsi que les identités
trigonométriques ont d'abord été obtenues. La démonstration classique introduit un point
auxiliaire ; le trouver est un rite de passage.

\begin{refs}
\item Contexte : Théorème de Ptolémée --- Wikipédia (\url{https://fr.wikipedia.org/wiki/Th%C3%A9or%C3%A8me_de_Ptol%C3%A9m%C3%A9e})
\item Lecture : H. S. M. Coxeter, \textit{Introduction to Geometry}, ch. 1
\end{refs}

\bigskip

\begin{problem}[{La droite d'Euler --- Lycée}]
\textbf{GEO-03.} \textbf{Problème.} Dans un triangle non équilatéral quelconque, soient $ O $ le centre du
cercle circonscrit (intersection des médiatrices), $ G $ le centre de gravité (intersection
des médianes) et $ H $ l'orthocentre (intersection des hauteurs).
\begin{enumerate}
\item[(a)] Démontrez que chacun de ces trois points existe bel et bien --- c'est-à-dire que les triplets
de droites concernés sont concourants.
\item[(b)] Démontrez que $ O, G, H $ sont alignés et que $ OG : GH = 1 : 2 $.
\end{enumerate}

\end{problem}

\noindent Euler a démontré cet alignement en 1765, deux millénaires après Euclide :
rappel saisissant que la géométrie élémentaire du triangle n'était pas close. La démonstration
moderne, élégante, repose sur une homothétie bien choisie de centre $ G $ ; découvrir cette
application enseigne davantage sur les transformations qu'un chapitre de théorie. La droite
porte aujourd'hui le nom d'Euler et passe par plusieurs autres centres remarquables.

\begin{refs}
\item Contexte : Droite d'Euler --- Wikipédia (\url{https://fr.wikipedia.org/wiki/Droite_d%27Euler})
\item Lecture : H. S. M. Coxeter, \textit{Introduction to Geometry}, ch. 1
\end{refs}

\bigskip

\begin{problem}[{Le cercle des neuf points --- Lycée}]
\textbf{GEO-04.} \textbf{Problème.} Considérez un triangle quelconque, d'orthocentre $ H $ et de centre du
cercle circonscrit $ O $.
\begin{enumerate}
\item[(a)] Démontrez que les neuf points suivants sont sur un même cercle : les milieux des trois
côtés, les pieds des trois hauteurs, et les milieux des trois segments joignant les sommets
à $ H $.
\item[(b)] Démontrez de plus que le centre de ce cercle est le milieu de $ OH $ (sur la droite
d'Euler de GEO-03) et que son rayon vaut la moitié du rayon du cercle circonscrit.
\end{enumerate}

\end{problem}

\noindent Le théorème complet a été publié par Karl Feuerbach en 1822, après des
versions partielles dues à Brianchon et Poncelet ; Feuerbach a ensuite montré que ce cercle
est tangent au cercle inscrit et aux trois cercles exinscrits, l'un des plus beaux résultats
de la géométrie classique. Une démonstration propre compose l'idée d'homothétie de GEO-03 avec
le théorème de l'angle inscrit --- la récompense d'avoir traité honnêtement les deux problèmes
précédents.

\begin{refs}
\item Contexte : Cercle d'Euler (cercle des neuf points) --- Wikipédia (\url{https://fr.wikipedia.org/wiki/Cercle_d%27Euler})
\item Lecture : H. S. M. Coxeter, \textit{Introduction to Geometry}, ch. 1
\end{refs}

\bigskip

\begin{problem}[{Isopérimétrie, premiers pas --- Lycée}]
\textbf{GEO-05.} \textbf{Problème.} Le but à l'horizon est le théorème isopérimétrique : parmi toutes les
courbes fermées de longueur donnée $ L $, le cercle enferme la plus grande aire, à savoir
$ L^2/4\pi $. Montez-y par étapes.
\begin{enumerate}
\item[(a)] Parmi tous les rectangles de périmètre donné, démontrez que le carré a la plus grande
aire.
\item[(b)] Parmi tous les triangles de base donnée et de périmètre donné, démontrez que l'isocèle a
la plus grande aire, et déduisez-en que parmi tous les triangles de périmètre donné c'est
l'équilatéral qui maximise l'aire.
\item[(c)] Démontrez qu'une région d'aire maximale à périmètre fixé est nécessairement convexe.
\item[(d)] Expliquez précisément ce qui manque encore à (a)--(c) si l'on veut le théorème
isopérimétrique complet.
\end{enumerate}

\end{problem}

\noindent Le problème isopérimétrique est antique --- la tradition grecque le rattache
à la reine Didon, à qui l'on accorda autant de terre qu'elle pourrait en entourer avec une
peau de bœuf découpée en lanières. Zénodore (IIe siècle av. J.-C.) comparait les
polygones ; Jakob Steiner (1838) donna de superbes arguments de symétrisation, dont l'astuce
de réflexion qui résout (c) --- mais l'argument de Steiner supposait, on le sait bien, l'existence
d'un maximiseur, lacune qui ne fut comblée que par Weierstrass. La question (d) vous demande de
trouver cette lacune vous-même : c'est une première leçon sur l'importance des démonstrations
d'existence.

\begin{refs}
\item Contexte : Théorème isopérimétrique --- Wikipédia (\url{https://fr.wikipedia.org/wiki/Th%C3%A9or%C3%A8me_isop%C3%A9rim%C3%A9trique})
\item Lecture : H. S. M. Coxeter, \textit{Introduction to Geometry}
\end{refs}

\bigskip

\begin{problem}[{La puissance d'un point --- Lycée, short}]
\textbf{GEO-23.} \textbf{Problème.} Soit $ P $ un point extérieur à un cercle, et soit une droite passant par $ P $ qui coupe le cercle en $ A $ et $ B $. Démontrez que le produit $ PA \cdot PB $ est le même pour toute droite de ce type passant par $ P $.
\end{problem}

\noindent Cet invariant est la puissance du point, et il transforme une famille de configurations apparemment sans lien en un seul nombre attaché à $ P $ seul. C'est le moteur de l'axe radical et de la plupart des problèmes d'olympiade mettant en jeu deux cercles, et la démonstration se réduit à une application des triangles semblables.

\begin{refs}
\item Contexte : Puissance d'un point par rapport à un cercle --- Wikipédia (\url{https://fr.wikipedia.org/wiki/Puissance_d%27un_point_par_rapport_%C3%A0_un_cercle})
\end{refs}

\bigskip

\begin{problem}[{Compter les points entiers dans un polygone --- Lycée, short}]
\textbf{GEO-24.} \textbf{Problème.} Vérifiez le théorème de Pick --- un polygone simple dont les sommets sont sur le réseau des points entiers a pour aire $ I + B/2 - 1 $, où $ I $ est le nombre de points entiers intérieurs et $ B $ le nombre de ceux du bord --- sur cinq polygones de votre invention, dont un non convexe. Démontrez-le ensuite pour les rectangles à côtés parallèles aux axes.
\end{problem}

\noindent Le théorème de Pick, démontré en 1899, relie une quantité continue (l'aire) à un dénombrement purement discret : c'est le genre de pont qui donnera plus tard toute la géométrie des points d'un réseau. Le vérifier sur un exemple non convexe est important : l'intuition de la plupart des gens suppose discrètement la convexité, alors que le théorème n'en a pas besoin.

\begin{refs}
\item Contexte : Théorème de Pick --- Wikipédia (\url{https://fr.wikipedia.org/wiki/Th%C3%A9or%C3%A8me_de_Pick})
\end{refs}

\bigskip

\begin{problem}[{Le théorème des trisectrices de Morley --- Lycée}]
\textbf{GEO-31.} \textbf{Problème.} Dans un triangle $ ABC $, partagez chacun des trois angles intérieurs en trois
parts égales. Pour chaque côté, prenez les deux trisectrices les plus proches de lui --- une issue de
chaque extrémité de ce côté --- et soient $ P_A, P_B, P_C $ les points où ces couples se rencontrent,
$ P_A $ étant celui associé au côté $ BC $.
\begin{enumerate}
\item[(a)] Démontrez que chacun de ces couples de trisectrices se rencontre effectivement à l'intérieur du
triangle, de sorte que les trois points existent.
\item[(b)] Démontrez le théorème de Morley : le triangle $ P_A P_B P_C $ est équilatéral, quel que soit le
triangle de départ. (Une voie trigonométrique utilise la loi des sinus jointe à l'identité
$ \sin 3\theta = 4 \sin\theta \, \sin(60^\circ + \theta) \, \sin(60^\circ - \theta) $.)
\item[(c)] Démontrez que le côté du triangle de Morley a pour longueur
\[ 8R \, \sin\tfrac{A}{3} \, \sin\tfrac{B}{3} \, \sin\tfrac{C}{3} , \]
où $ R $ est le rayon du cercle circonscrit à $ ABC $ et $ A, B, C $ ses angles.
\item[(d)] Autorisez également les trisectrices des angles \textit{extérieurs}. Exhibez au moins un autre
triangle équilatéral obtenu de cette façon --- il y a en tout dix-huit triangles de Morley --- et dites
lesquels de vos arguments de (b) se transposent sans changement.
\end{enumerate}

\end{problem}

\noindent Frank Morley, un Anglais qui fit toute sa carrière à Haverford puis à Johns Hopkins,
remarqua ce fait vers 1899 en étudiant la géométrie des cardioïdes tangentes aux côtés d'un triangle ;
il en parla à des amis et le résultat circula de bouche à oreille bien avant sa publication, la
première démonstration imprimée étant due à M. T. Naraniengar en 1909. Ce retard n'a rien d'un
hasard : la trisection de l'angle est impossible à la règle et au compas, si bien que la configuration
ne peut être tracée par les moyens classiques et qu'aucun Grec n'aurait pu tomber dessus. Le théorème
attire encore de nouvelles démonstrations --- Alain Connes en a publié une, algébrique, en 1998, qui
vaut sur tout corps de caractéristique différente de deux et de trois.

\begin{refs}
\item Contexte : Théorème de Morley --- Wikipédia (\url{https://fr.wikipedia.org/wiki/Th%C3%A9or%C3%A8me_de_Morley})
\item Contexte : Trisection de l'angle --- Wikipédia (\url{https://fr.wikipedia.org/wiki/Trisection_de_l%27angle})
\item Lecture : H. S. M. Coxeter et S. L. Greitzer, \textit{Geometry Revisited}, ch. 1
\item Lecture : H. S. M. Coxeter, \textit{Introduction to Geometry}, ch. 1
\end{refs}

\bigskip

\begin{problem}[{Deux bissectrices de même longueur --- Lycée, short}]
\textbf{GEO-32.} \textbf{Problème.} Démontrez le théorème de Steiner--Lehmus : si les bissectrices
intérieures issues de deux sommets d'un triangle ont la même longueur, alors le triangle est
isocèle. (La réciproque est une conséquence immédiate de la symétrie ; toute la difficulté est dans
ce sens-ci.)
\end{problem}

\noindent C. L. Lehmus demanda à Sturm, dans une lettre de 1840, une démonstration purement
géométrique ; Sturm fit circuler la question, et Jakob Steiner fut parmi les premiers à répondre ---
depuis, le problème a suscité bien plus d'une centaine de démonstrations publiées. Presque toutes
raisonnent par l'absurde, en prouvant plutôt qu'un triangle non isocèle a des bissectrices de
longueurs différentes, et la possibilité d'une démonstration véritablement \og{} directe \fg{} est devenue
une petite controverse logique à part entière : John Conway soutenait qu'il n'en existe aucune,
Victor Pambuccian a démontré ensuite que des démonstrations directes doivent exister en logique
classique comme en logique intuitionniste, et l'on en a depuis écrit une vérifiée par
machine.

\begin{refs}
\item Contexte : Théorème de Steiner-Lehmus --- Wikipédia (\url{https://fr.wikipedia.org/wiki/Th%C3%A9or%C3%A8me_de_Steiner-Lehmus})
\item Lecture : H. S. M. Coxeter et S. L. Greitzer, \textit{Geometry Revisited}, ch. 1
\end{refs}

\bigskip


\section*{Licence}

\begin{problem}[{Le postulat des parallèles et l'axiome de Playfair --- Licence}]
\textbf{GEO-06.} \textbf{Problème.} Travaillez à partir des quatre premiers postulats d'Euclide --- ou, pour être
tout à fait rigoureux, à partir des axiomes d'incidence, d'ordre et de congruence de Hilbert
(la \og{} géométrie neutre \fg{}).
\begin{enumerate}
\item[(a)] Démontrez que le cinquième postulat d'Euclide équivaut à l'axiome de Playfair : par un
point extérieur à une droite $ \ell $ il passe au plus une droite qui ne rencontre pas
$ \ell $.
\item[(b)] Démontrez que les deux équivalent à l'énoncé \og{} il existe un triangle dont la somme des
angles vaut exactement $ 180^\circ $ \fg{}.
\end{enumerate}

\end{problem}

\noindent Pendant deux mille ans, les géomètres ont tenté de \textit{démontrer} le cinquième
postulat à partir des autres --- Proclus, Omar Khayyam, Saccheri (dont l'ouvrage de 1733 \og{} vengeant
Euclide \fg{} déroule des pages de géométrie hyperbolique correcte tout en la croyant absurde),
Lambert, Legendre. Les équivalences ci-dessus expliquent pourquoi chaque tentative supposait
secrètement la conclusion. La résolution --- le postulat est indépendant, et le nier donne une
géométrie cohérente --- est venue avec Lobatchevski (1829) et János Bolyai (1832), Gauss en étant
convaincu en privé plus tôt. Les \textit{Grundlagen} de Hilbert (1899) ont enfin fait de
\og{} travailler à partir des axiomes \fg{} une activité précise.

\begin{refs}
\item Contexte : Axiome des parallèles --- Wikipédia (\url{https://fr.wikipedia.org/wiki/Axiome_des_parall%C3%A8les})
\item Contexte : Axiome de Playfair --- Wikipédia (en anglais) (\url{https://en.wikipedia.org/wiki/Playfair%27s_axiom})
\item Lecture : R. Hartshorne, \textit{Geometry: Euclid and Beyond}, ch. 1--4
\end{refs}

\bigskip

\begin{problem}[{Les géodésiques du plan hyperbolique --- Licence}]
\textbf{GEO-07.} \textbf{Problème.} Le demi-plan supérieur $ \mathbb{H} = \{(x,y) : y > 0\} $, muni de la
longueur de courbe définie comme la longueur euclidienne usuelle pondérée en chaque point par
$ 1/y $, est un modèle du plan hyperbolique.
\begin{enumerate}
\item[(a)] Démontrez que les géodésiques (les courbes localement les plus courtes) sont exactement les
demi-droites verticales et les demi-cercles rencontrant l'axe des $ x $ à angle droit.
\item[(b)] Concluez que par tout point $ P $ hors d'une géodésique $ g $ il passe une infinité de
géodésiques disjointes de $ g $ --- l'axiome de Playfair (GEO-06) est donc en défaut, tandis que
les quatre premiers postulats d'Euclide, convenablement interprétés, restent vrais.
\end{enumerate}

\end{problem}

\noindent Les modèles sont ce qui a rendu la géométrie non euclidienne respectable :
Beltrami (1868), puis Klein et Poincaré, ont montré que toute contradiction dans la géométrie
hyperbolique serait déjà une contradiction dans la géométrie euclidienne, réglant la question
bimillénaire de GEO-06. Les modèles du demi-plan et du disque restent aujourd'hui les images de
l'espace hyperbolique dont se sert le géomètre au travail --- dans le modèle du disque, les gravures
\textit{Circle Limit} d'Escher sont littéralement des dessins d'un pavage hyperbolique.

\begin{refs}
\item Contexte : Demi-plan de Poincaré --- Wikipédia (\url{https://fr.wikipedia.org/wiki/Demi-plan_de_Poincar%C3%A9})
\item Contexte : Géométrie hyperbolique --- Wikipédia (\url{https://fr.wikipedia.org/wiki/G%C3%A9om%C3%A9trie_hyperbolique})
\item Lecture : R. Hartshorne, \textit{Geometry: Euclid and Beyond}, ch. 7
\end{refs}

\bigskip

\begin{problem}[{Le défaut angulaire est une aire --- Licence}]
\textbf{GEO-08.} \textbf{Problème.} Travaillez dans le plan hyperbolique, avec le modèle du demi-plan de GEO-07,
dont l'élément d'aire est l'élément euclidien multiplié par $ 1/y^2 $.
\begin{enumerate}
\item[(a)] Démontrez que tout triangle a une somme d'angles strictement inférieure à $ \pi $.
\item[(b)] Démontrez que l'aire d'un triangle d'angles $ \alpha, \beta, \gamma $ est exactement son
défaut angulaire
\[ \pi - (\alpha + \beta + \gamma). \]
Commencez par les triangles idéaux --- dont les trois sommets sont sur le bord --- qui ont tous
leurs angles $ 0 $ et, comme vous le montrerez, une aire $ \pi $.
\item[(c)] Concluez que les triangles hyperboliques ont une aire inférieure à $ \pi $, quelle que soit
la longueur de leurs côtés.
\end{enumerate}

\end{problem}

\noindent Lambert observa dans les années 1760 que, si le postulat des parallèles tombe,
le défaut angulaire doit se comporter comme une aire --- et nota l'étrange analogie avec l'excès
sphérique (GEO-09), spéculant sur une \og{} sphère de rayon imaginaire \fg{}. Que les triangles ne
puissent pas avoir une grande aire est l'un des premiers faits véritablement étranges que
rencontre l'étudiant en géométrie hyperbolique, et c'est le cas de courbure $ -1 $ du théorème
de Gauss--Bonnet (GEO-14).

\begin{refs}
\item Contexte : Géométrie hyperbolique --- Wikipédia (\url{https://fr.wikipedia.org/wiki/G%C3%A9om%C3%A9trie_hyperbolique})
\item Lecture : R. Hartshorne, \textit{Geometry: Euclid and Beyond}, ch. 7
\item Lecture : H. S. M. Coxeter, \textit{Introduction to Geometry}, ch. 16
\end{refs}

\bigskip

\begin{problem}[{Le théorème de Girard sur la sphère --- Licence}]
\textbf{GEO-09.} \textbf{Problème.} Un fuseau sur une sphère est une région bornée par deux demi-grands cercles de
mêmes extrémités ; un triangle sphérique est borné par trois arcs de grands cercles.
\begin{enumerate}
\item[(a)] Démontrez le théorème de Girard : sur une sphère de rayon $ R $, un triangle d'angles
$ \alpha, \beta, \gamma $ a pour aire
\[ R^2\,(\alpha + \beta + \gamma - \pi). \]
(L'indication est dans l'énoncé : calculez d'abord l'aire d'un fuseau, puis recouvrez la sphère
par les fuseaux engendrés par les trois côtés du triangle.)
\item[(b)] Déduisez-en que la somme des angles d'un triangle sphérique est toujours supérieure à
$ 180^\circ $.
\item[(c)] Déduisez-en que sur la sphère il n'existe pas de triangles semblables mais non égaux.
\end{enumerate}

\end{problem}

\noindent Thomas Harriot démontra ce résultat en 1603 mais ne le publia jamais ; Albert
Girard l'imprima en 1629. C'est le plus ancien théorème d'une géométrie non euclidienne, découvert
des siècles avant que quiconque ne le formule ainsi --- les navigateurs avaient simplement besoin de
trigonométrie sphérique. La quantité $ \alpha + \beta + \gamma - \pi $ est l'\og{} excès sphérique \fg{},
miroir exact du défaut hyperbolique (GEO-08), et les mesures géodésiques de Gauss sur un grand
triangle reliant trois sommets allemands n'étaient assez précises que pour confirmer que l'espace
physique est très proche d'être plat.

\begin{refs}
\item Contexte : Trigonométrie sphérique --- Wikipédia (\url{https://fr.wikipedia.org/wiki/Trigonom%C3%A9trie_sph%C3%A9rique})
\item Contexte : Géométrie sphérique --- Wikipédia (\url{https://fr.wikipedia.org/wiki/G%C3%A9om%C3%A9trie_sph%C3%A9rique})
\item Lecture : H. S. M. Coxeter, \textit{Introduction to Geometry}
\end{refs}

\bigskip

\begin{problem}[{La courbure détermine la courbe --- Licence}]
\textbf{GEO-10.} \textbf{Problème.} Soit une courbe plane paramétrée par l'abscisse curviligne $ s $, et soit
$ \kappa(s) $ sa courbure algébrique --- la vitesse à laquelle tourne son angle de direction.
\begin{enumerate}
\item[(a)] Démontrez le théorème fondamental des courbes planes : étant donné une fonction continue
$ \kappa(s) $ quelconque, il existe une courbe de vitesse unité ayant cette courbure, unique à
un déplacement du plan près.
\item[(b)] Déduisez-en que les seules courbes planes de courbure constante sont les droites et les
cercles.
\item[(c)] Pour une courbe de l'espace, définissez la courbure et la torsion via le repère de
Frenet--Serret, puis énoncez (ou démontrez) le théorème analogue.
\end{enumerate}

\end{problem}

\noindent C'est la porte d'entrée de la géométrie différentielle : des invariants locaux
(courbure, torsion) déterminent la forme globale. Le cas plan est un agréable exercice d'équations
différentielles ordinaires ; le cas de l'espace, systématisé par Frenet (1847) et Serret (1851),
est le modèle que Gauss et Riemann ont généralisé aux surfaces et au-delà. La morale --- la géométrie
est stockée dans les dérivées --- c'est tout le sujet en une phrase.

\begin{refs}
\item Contexte : Courbure --- Wikipédia (\url{https://fr.wikipedia.org/wiki/Courbure})
\item Contexte : Repère de Frenet --- Wikipédia (\url{https://fr.wikipedia.org/wiki/Rep%C3%A8re_de_Frenet})
\item Lecture : M. do Carmo, \textit{Differential Geometry of Curves and Surfaces}, ch. 1
\end{refs}

\bigskip

\begin{problem}[{Géodésiques : grands cercles et relation de Clairaut --- Licence}]
\textbf{GEO-11.} \textbf{Problème.} Une géodésique sur une surface est une courbe localement la
plus courte, ce qui équivaut à dire que son accélération est partout normale à la surface.
(a) Démontrez que les géodésiques de la sphère ronde sont exactement les grands cercles. (Une voie
élégante : montrez que l'ensemble des points fixes d'une isométrie qui fixe le point de départ et
la direction d'une géodésique contient nécessairement cette géodésique.) (b) Sur une surface de
révolution, démontrez la relation de Clairaut : le long de toute géodésique, la quantité
$ r \cos\theta $ est constante, où $ r $ est la distance à l'axe et $ \theta $ l'angle que
fait la géodésique avec le parallèle local. Servez-vous-en pour décrire qualitativement comment
les géodésiques d'une surface en forme de trompette rebondissent entre les parallèles.
\end{problem}

\noindent Les géodésiques répondent à la question du marin --- quelle est la route la plus
courte ? --- et expliquent pourquoi les vols long-courriers décrivent un arc vers les pôles sur une
carte plate. Clairaut publia sa relation en 1733, avant même que la notion générale de courbure
géodésique n'existe ; c'est la conservation du moment cinétique déguisée, exemple précoce du
principe symétrie--conservation que Noether rendra plus tard précis. Les géodésiques sont aussi les
trajectoires de la lumière et des particules en chute libre en relativité générale.

\begin{refs}
\item Contexte : Géodésique --- Wikipédia (\url{https://fr.wikipedia.org/wiki/G%C3%A9od%C3%A9sique})
\item Lecture : M. do Carmo, \textit{Differential Geometry of Curves and Surfaces}, ch. 4
\item Cours : MIT OCW 18.950 Differential Geometry (\url{https://ocw.mit.edu/})
\end{refs}

\bigskip

\begin{problem}[{Les sphères de Dandelin : les coniques à partir du cône --- Licence}]
\textbf{GEO-12.} \textbf{Problème.} Un plan coupe une nappe d'un cône de révolution suivant une
courbe bornée $ C $, rencontrant chaque génératrice une fois. Inscrivez dans le cône deux sphères
tangentes au plan de coupe, une de chaque côté, touchant le plan aux points $ F_1 $ et
$ F_2 $. Démontrez que pour tout point $ P $ de $ C $, la somme $ PF_1 + PF_2 $ est
constante --- de sorte que $ C $ est une ellipse de foyers $ F_1, F_2 $. Adaptez l'argument pour
obtenir la propriété foyer--directrice, et décrivez ce qui change pour la parabole et
l'hyperbole.
\end{problem}

\noindent Les Grecs --- Ménechme, Euclide et surtout Apollonios de Perga (\textit{Les Coniques},
vers 200 av. J.-C.) --- définissaient les coniques comme sections d'un cône mais en tiraient les
propriétés focales par de longs arguments synthétiques. Les deux sphères de Germinal Dandelin
(1822) réduisent la dérivation à quelques lignes de géométrie pure : c'est une sérieuse candidate
au titre de plus élégante démonstration du domaine. Les coniques sont aussi la porte d'entrée de la
géométrie algébrique : ce sont exactement les courbes planes de degré 2, premier cas des questions
de classification de GEO-16.

\begin{refs}
\item Contexte : Théorème de Dandelin (sphères de Dandelin) --- Wikipédia (\url{https://fr.wikipedia.org/wiki/Th%C3%A9or%C3%A8me_de_Dandelin})
\item Contexte : Conique --- Wikipédia (\url{https://fr.wikipedia.org/wiki/Conique})
\item Lecture : H. S. M. Coxeter, \textit{Introduction to Geometry}
\end{refs}

\bigskip

\begin{problem}[{Toute isométrie du plan est l'une de quatre choses --- Licence, short}]
\textbf{GEO-25.} \textbf{Problème.} Démontrez que toute isométrie du plan euclidien est une translation, une rotation, une réflexion ou une réflexion glissée, et que toute isométrie est composée d'au plus trois réflexions.
\end{problem}

\noindent Voilà une classification complète obtenue avec presque aucun outillage, et elle explique pourquoi les dix-sept groupes de papier peint sont ce qu'ils sont. L'énoncé \og{} trois réflexions \fg{} est la forme mémorable : les réflexions engendrent tout, si bien que le groupe des symétries du plan tout entier se construit à partir d'un seul type de mouvement.

\begin{refs}
\item Contexte : Isométrie du plan euclidien --- Wikipédia (en anglais) (\url{https://en.wikipedia.org/wiki/Euclidean_plane_isometry})
\end{refs}

\bigskip

\begin{problem}[{Helly dans le plan --- Licence, short}]
\textbf{GEO-26.} \textbf{Problème.} Démontrez le théorème de Helly dans le plan : si une famille finie de convexes de $ \mathbb{R}^2 $ est telle que trois quelconques d'entre eux ont un point commun, alors tous ont un point commun. Montrez par un exemple que \og{} trois \fg{} ne peut pas être affaibli en \og{} deux \fg{}.
\end{problem}

\noindent Le théorème de Helly, de 1913, est le prototype de toute une famille de résultats affirmant que la convexité permet à un accord local de forcer un accord global. Le contre-exemple est aussi instructif que la démonstration --- trois droites en position générale se coupent deux à deux sans avoir de point commun --- car il montre que la constante, qui dépend de la dimension, est optimale.

\begin{refs}
\item Contexte : Théorème de Helly --- Wikipédia (\url{https://fr.wikipedia.org/wiki/Th%C3%A9or%C3%A8me_de_Helly})
\end{refs}

\bigskip

\begin{problem}[{Naviguer sur une Terre ronde --- Licence}]
\textbf{GEO-33.} \textbf{Problème.} Modélisez la Terre par une sphère de rayon $ R = 6371 $ km. Le point de
latitude $ \varphi $ et de longitude $ \lambda $ est le vecteur unitaire
$ (\cos\varphi\cos\lambda, \ \cos\varphi\sin\lambda, \ \sin\varphi) $, et la distance entre deux
points vaut $ R $ fois l'angle entre leurs vecteurs.
\begin{enumerate}
\item[(a)] Démontrez la loi des cosinus sphérique : dans un triangle sphérique de côtés $ a, b, c $
(mesurés comme angles au centre) et où $ C $ est l'angle au sommet opposé au côté $ c $,
\[ \cos c = \cos a \cos b + \sin a \sin b \cos C . \]
Obtenez-la à partir des produits scalaire et vectoriel des trois vecteurs unitaires.
\item[(b)] Déduisez-en la formule de la distance orthodromique : pour les points $ (\varphi_1, \lambda_1) $
et $ (\varphi_2, \lambda_2) $,
\[ \frac{d}{R} = \arccos\Big( \sin\varphi_1 \sin\varphi_2 + \cos\varphi_1 \cos\varphi_2
\cos(\lambda_2 - \lambda_1) \Big) . \]
\item[(c)] Calculez la distance orthodromique de Londres ($ 51{,}5^\circ $ N, $ 0{,}13^\circ $ O) à
New York ($ 40{,}71^\circ $ N, $ 74{,}01^\circ $ O), et comparez-la à la longueur de la
\textit{loxodromie} --- la route à cap constant --- entre ces deux mêmes points. Expliquez
géométriquement pourquoi la route la plus courte commence par un cap nettement au nord de
l'ouest.
\item[(d)] Montrez que la formule de (b) perd en précision pour des points proches, l'arccosinus étant
mal conditionné au voisinage de $ 1 $, et démontrez que l'identité de haversine
\[ \sin^2\!\Big(\frac{d}{2R}\Big) = \sin^2\!\Big(\frac{\varphi_2 - \varphi_1}{2}\Big)
+ \cos\varphi_1 \cos\varphi_2 \, \sin^2\!\Big(\frac{\lambda_2 - \lambda_1}{2}\Big) \]
lui est équivalente tout en étant numériquement sûre.
\end{enumerate}

\end{problem}

\noindent La trigonométrie sphérique est plus ancienne que la géométrie analytique plane :
la loi des cosinus pour les côtés était en usage chez les astronomes du monde islamique médiéval
calculant la direction de La Mecque, et toute la discipline existait pour répondre à des questions
sur le ciel et sur la mer. La loxodromie fut décrite par Pedro Nunes en 1537, et la projection de
Mercator de 1569 fut conçue précisément pour la tracer comme une droite --- une carte qui échange les
routes les plus courtes contre des routes gouvernables. La forme haversine de (d) fut mise en table
pour les navigateurs au début du XIXe siècle et nommée par James Inman en 1835 ; c'est
encore la formule employée par la bibliothèque cartographique de votre téléphone, exactement pour
la raison qu'on vous demande de démontrer.

\begin{refs}
\item Contexte : Loi des cosinus sphérique --- Wikipédia (en anglais) (\url{https://en.wikipedia.org/wiki/Spherical_law_of_cosines}), Distance du grand cercle --- Wikipédia (\url{https://fr.wikipedia.org/wiki/Distance_du_grand_cercle})
\item Contexte : Formule de haversine --- Wikipédia (\url{https://fr.wikipedia.org/wiki/Formule_de_haversine}), Loxodromie --- Wikipédia (\url{https://fr.wikipedia.org/wiki/Loxodromie})
\item Lecture : G. Van Brummelen, \textit{Heavenly Mathematics: The Forgotten Art of Spherical Trigonometry} (Princeton, 2013)
\item Lecture : H. S. M. Coxeter, \textit{Introduction to Geometry}, ch. 6
\end{refs}

\bigskip

\begin{problem}[{L'inversion, et le porisme de Steiner --- Licence}]
\textbf{GEO-34.} \textbf{Problème.} L'inversion de centre $ O $ et de rayon $ r $ est l'application du plan
privé de $ O $ dans lui-même donnée par
\[ P \ \longmapsto \ O + r^2 \, \frac{P - O}{|P - O|^2} . \]
\begin{enumerate}
\item[(a)] Démontrez que l'inversion est une involution, qu'elle envoie les cercles ne passant pas par
$ O $ sur des cercles, les cercles passant par $ O $ sur des droites évitant $ O $, et les
droites passant par $ O $ sur elles-mêmes --- de sorte que \og{} cercles et droites \fg{} forment une seule
famille, permutée par les inversions.
\item[(b)] Démontrez que l'inversion conserve la valeur absolue de l'angle entre deux courbes, tout en en
renversant l'orientation.
\item[(c)] Soient $ \alpha $ et $ \beta $ deux cercles, $ \beta $ intérieur à $ \alpha $ et les
deux non concentriques. Démontrez qu'il existe une inversion les transformant en deux cercles
concentriques.
\item[(d)] Une \textit{chaîne de Steiner} est une suite cyclique de cercles, chacun tangent au suivant, tous
intérieurement tangents à $ \alpha $ et extérieurement tangents à $ \beta $. À l'aide de (c),
démontrez le porisme de Steiner : s'il existe une chaîne fermée de $ n $ cercles pour le couple
$ (\alpha, \beta) $, alors une chaîne fermée de $ n $ cercles part de \textit{n'importe quel}
cercle tangent aux deux --- la chaîne se referme toujours, quelle que soit la position de départ.
Démontrez aussi que les centres des cercles de la chaîne sont sur une même ellipse dont les foyers
sont les centres de $ \alpha $ et de $ \beta $.
\end{enumerate}

\end{problem}

\noindent Jakob Steiner découvrit ce porisme dans les années 1820, dans l'élan de travaux
qui produisit aussi les arguments de symétrisation de GEO-05. Le mot \og{} porisme \fg{} vient d'un livre
perdu d'Euclide de ce nom et décrit exactement ce type d'énoncé : non pas \og{} voici l'objet \fg{} mais
\og{} s'il en existe un, il en existe une infinité \fg{}. L'inversion est l'outil qui rend la chose évidente
au lieu de miraculeuse --- vous transportez la configuration difficile en une configuration facile,
vous la résolvez par symétrie, et vous revenez --- et la même astuce donne la démonstration standard
de l'inégalité de Ptolémée de GEO-02(c) et du théorème de tangence de Feuerbach derrière GEO-04.
C'est aussi la première transformation véritablement non euclidienne que rencontrent la plupart des
étudiants : les angles survivent, les distances et l'alignement non.

\begin{refs}
\item Contexte : Inversion géométrique --- Wikipédia (\url{https://fr.wikipedia.org/wiki/Inversion_g%C3%A9om%C3%A9trique})
\item Contexte : Chaîne de Steiner --- Wikipédia (\url{https://fr.wikipedia.org/wiki/Cha%C3%AEne_de_Steiner})
\item Lecture : H. S. M. Coxeter et S. L. Greitzer, \textit{Geometry Revisited}, ch. 5
\item Lecture : H. S. M. Coxeter, \textit{Introduction to Geometry}, ch. 6
\end{refs}

\bigskip

\begin{problem}[{L'inégalité d'Erdős--Mordell --- Licence, short}]
\textbf{GEO-35.} \textbf{Problème.} Soit $ P $ un point intérieur à un triangle $ ABC $, et
soient $ p_a, p_b, p_c $ les distances de $ P $ aux trois côtés. Démontrez que
\[ PA + PB + PC \ \ge \ 2\,(p_a + p_b + p_c) , \]
et déterminez exactement le cas d'égalité.
\end{problem}

\noindent Paul Erdős proposa cet énoncé comme problème dans l'\textit{American Mathematical
Monthly} en 1935 sans en avoir de démonstration ; Louis Mordell et D. F. Barrow en fournirent
une en 1937, au prix d'une trigonométrie assez lourde pour que trouver un argument élémentaire
devienne un sport, et l'on dispose aujourd'hui de démonstrations courtes par triangles semblables
ou par l'inégalité arithmético-géométrique. La constante $ 2 $ est ce qui rend l'inégalité fine
et délicate : le cas d'égalité fixe une seule configuration, et l'inégalité de Barrow renforce
l'énoncé en remplaçant les distances aux côtés par des segments plus longs.

\begin{refs}
\item Contexte : Théorème d'Erdős-Mordell --- Wikipédia (\url{https://fr.wikipedia.org/wiki/Th%C3%A9or%C3%A8me_d%27Erd%C5%91s-Mordell})
\item Lecture : H. S. M. Coxeter et S. L. Greitzer, \textit{Geometry Revisited}
\end{refs}

\bigskip


\section*{Master}

\begin{problem}[{Le Theorema Egregium, ou pourquoi toute carte de la Terre ment --- Master}]
\textbf{GEO-13.} \textbf{Problème.} (Problème de compréhension.) La courbure de Gauss $ K $
d'une surface de l'espace est le produit de ses deux courbures principales. Travaillez une
démonstration complète du Theorema Egregium de Gauss : $ K $ est intrinsèque --- calculable à
partir de mesures de longueur faites au sein de la surface, donc invariante par flexion sans
étirement. Rédigez ensuite, en une page autonome, l'argument du corollaire : aucune carte d'une
région de la surface terrestre, si petite soit-elle, ne peut représenter toutes les distances à une
même échelle. Enfin, conciliez cela avec la pratique cartographique : identifiez une projection qui
conserve tous les angles et une autre qui conserve toutes les aires, et démontrez qu'aucune
projection ne peut faire les deux.
\end{problem}

\noindent Gauss publia le théorème dans ses \textit{Disquisitiones generales circa superficies
curvas} (1827) et le baptisa lui-même \og{} egregium \fg{} --- remarquable : la courbure était définie en
faisant appel à l'espace ambiant, et il se trouve qu'elle n'en a pas besoin. Le corollaire donne
son statut professionnel à un fait que tout usager d'atlas pressent (le Groenland n'est pas si
grand ; Mercator conserve les angles au prix brutal des aires). La leçon de 1854 de Riemann prit le
point de vue intrinsèque comme définition, créant la géométrie riemannienne, langue dans laquelle
la relativité générale allait s'écrire.

\begin{refs}
\item Contexte : Theorema egregium --- Wikipédia (\url{https://fr.wikipedia.org/wiki/Theorema_egregium})
\item Contexte : Projection cartographique --- Wikipédia (\url{https://fr.wikipedia.org/wiki/Projection_cartographique})
\item Lecture : M. do Carmo, \textit{Differential Geometry of Curves and Surfaces}, ch. 4
\end{refs}

\bigskip

\begin{problem}[{Le théorème de Gauss--Bonnet --- Master}]
\textbf{GEO-14.} \textbf{Problème.} Admettez la formule de Gauss--Bonnet locale : pour un triangle géodésique
$ T $ sur une surface, $ \iint_T K \, dA = \alpha + \beta + \gamma - \pi $. Par
triangulation, démontrez le théorème global : pour une surface compacte orientable $ S $ de
genre $ g $,
\[ \iint_S K \, dA = 2\pi\,\chi(S) = 2\pi\,(2 - 2g). \]
Déduisez-en :
\begin{enumerate}
\item[(a)] que la courbure totale de toute surface de type sphère vaut $ 4\pi $ et celle de tout tore
$ 0 $, quelle que soit la manière dont on les déforme ;
\item[(b)] qu'une surface de genre au moins $ 2 $ doit porter quelque part de la courbure
négative ;
\item[(c)] que le théorème de Girard (GEO-09) et la formule du défaut hyperbolique (GEO-08) en sont les
cas particuliers à courbure constante.
\end{enumerate}

\end{problem}

\noindent Bonnet démontra la version à bord en 1848, en s'appuyant sur Gauss ; l'énoncé
pour les surfaces fermées est le prototype du thème le plus profond de la géométrie moderne, le
principe du local au global : la courbure est locale et librement déformable, et pourtant son total
est un invariant topologique, rigide comme un entier. Sa descendance --- Chern--Gauss--Bonnet en
dimension supérieure, Riemann--Roch, le théorème de l'indice d'Atiyah--Singer --- forme l'un des grands
récits des mathématiques du XXe siècle.

\begin{refs}
\item Contexte : Formule de Gauss-Bonnet --- Wikipédia (\url{https://fr.wikipedia.org/wiki/Formule_de_Gauss-Bonnet})
\item Lecture : M. do Carmo, \textit{Differential Geometry of Curves and Surfaces}, ch. 4
\item Cours : MIT OCW 18.950 Differential Geometry (\url{https://ocw.mit.edu/})
\end{refs}

\bigskip

\begin{problem}[{Surfaces minimales et films de savon --- Master}]
\textbf{GEO-15.} \textbf{Problème.} Une surface est minimale si sa courbure moyenne (la moyenne des courbures
principales) est identiquement nulle --- condition d'équilibre d'un film de savon, que la physique
contraint à minimiser l'aire localement.
\begin{enumerate}
\item[(a)] Démontrez qu'une surface de révolution est minimale si et seulement si c'est un plan ou une
caténoïde, la surface engendrée par la rotation du graphe de $ \cosh $.
\item[(b)] Démontrez que l'hélicoïde --- la surface en escalier en colimaçon --- est minimale.
\item[(c)] Démontrez qu'une surface minimale de $ \mathbb{R}^3 $ ne peut avoir de composante compacte,
de sorte que tout film de savon doit avoir un fil de bord.
\end{enumerate}

\end{problem}

\noindent Euler trouva la caténoïde en 1744, première surface minimale après le plan ; le
mémoire de Lagrange de 1762 fondant le calcul des variations posa le problème général de l'aire
minimale, et Meusnier l'interpréta par la courbure moyenne en 1776. Le physicien belge aveugle
Joseph Plateau (années 1870) fit des films de savon une branche expérimentale des mathématiques :
trempez n'importe quel cadre de fil de fer et une surface minimale apparaît en une seconde. Savoir
si les mathématiques pouvaient égaler l'expérience --- l'existence d'une surface d'aire minimale pour
toute courbe de bord --- est devenu le problème de Plateau, objet de GEO-22.

\begin{refs}
\item Contexte : Surface minimale --- Wikipédia (\url{https://fr.wikipedia.org/wiki/Surface_minimale})
\item Contexte : Caténoïde --- Wikipédia (\url{https://fr.wikipedia.org/wiki/Cat%C3%A9no%C3%AFde})
\item Lecture : M. do Carmo, \textit{Differential Geometry of Curves and Surfaces}, ch. 3
\end{refs}

\bigskip

\begin{problem}[{Le théorème de Bézout, mis à l'épreuve à la main --- Master}]
\textbf{GEO-16.} \textbf{Problème.} Travaillez dans le plan projectif complexe $ \mathbb{P}^2(\mathbb{C}) $.
\begin{enumerate}
\item[(a)] Démontrez qu'une droite rencontre une courbe de degré $ d $ ne la contenant pas en
exactement $ d $ points, comptés avec multiplicité.
\item[(b)] Démontrez que deux coniques sans composante commune se rencontrent en exactement $ 4 $
points, comptés avec multiplicité, et exhibez des configurations réalisant les schémas
d'intersection $ 1{+}1{+}1{+}1 $, $ 2{+}1{+}1 $, $ 2{+}2 $, $ 3{+}1 $ et $ 4 $.
\item[(c)] Énoncez le théorème de Bézout --- deux courbes de degrés $ m $ et $ n $ sans composante
commune se rencontrent en exactement $ mn $ points comptés avec multiplicité --- et montrez par des
exemples explicites que chacune des hypothèses (travailler sur $ \mathbb{C} $, travailler
projectivement, compter les multiplicités) est nécessaire.
\end{enumerate}

\end{problem}

\noindent Étienne Bézout énonça ce décompte en 1779 ; le rendre vrai a demandé un siècle
d'invention des bonnes notions --- points à l'infini (espace projectif), points complexes et
multiplicité d'intersection. Cette trajectoire, \og{} réparer le théorème en réparant les définitions \fg{},
est le geste fondateur de la géométrie algébrique, et ce problème vous invite à le revivre. Bézout
est aussi le moteur de GEO-17 : la loi de groupe sur une cubique repose sur le fait qu'une droite
la rencontre en exactement 3 points.

\begin{refs}
\item Contexte : Théorème de Bézout --- Wikipédia (\url{https://fr.wikipedia.org/wiki/Th%C3%A9or%C3%A8me_de_B%C3%A9zout})
\item Lecture : M. Reid, \textit{Undergraduate Algebraic Geometry}
\item Cours : MIT OCW 18.725 Algebraic Geometry (\url{https://ocw.mit.edu/})
\end{refs}

\bigskip

\begin{problem}[{La loi de groupe sur une cubique plane --- Master}]
\textbf{GEO-17.} \textbf{Problème.} Soit $ E $ une cubique non singulière de
$ \mathbb{P}^2(\mathbb{C}) $, et fixez un point $ O $ sur $ E $. Définissez l'addition par
cordes et tangentes : pour $ P, Q $ sur $ E $, soit $ P \ast Q $ la troisième intersection de
la droite $ PQ $ avec $ E $ (Bézout, GEO-16), et posez $ P + Q = O \ast (P \ast Q) $.
Démontrez que $ (E, +) $ est un groupe abélien d'élément neutre $ O $. Tout est routinier sauf
l'associativité --- démontrez-la, soit via le théorème de Cayley--Bacharach (une cubique passant par
8 des 9 points d'intersection de deux cubiques passe par le 9e), soit par un calcul
honnête que vous êtes prêt à défendre.
\end{problem}

\noindent La construction par cordes est implicite chez Diophante et explicite chez Fermat
et Newton comme machine à fabriquer de nouveaux points rationnels à partir d'anciens ; Poincaré
(1901) a dégagé la structure de groupe comme le bon objet d'étude. Les courbes elliptiques sont
aujourd'hui à un carrefour des mathématiques : théorème de Mordell, conjecture de
Birch--Swinnerton-Dyer (un problème du prix du millénaire --- voir la page des problèmes ouverts),
démonstration par Wiles du dernier théorème de Fermat, et la cryptographie sur courbes elliptiques
de votre téléphone. L'associativité est la seule étape vraiment difficile, et tout étudiant sérieux
devrait en souffrir une fois.

\begin{refs}
\item Contexte : Courbe elliptique --- Wikipédia (\url{https://fr.wikipedia.org/wiki/Courbe_elliptique})
\item Contexte : Théorème de Cayley--Bacharach --- Wikipédia (en anglais) (\url{https://en.wikipedia.org/wiki/Cayley%E2%80%93Bacharach_theorem})
\item Lecture : M. Reid, \textit{Undergraduate Algebraic Geometry}, ch. 2
\end{refs}

\bigskip

\begin{problem}[{Vingt-sept droites sur une surface cubique --- Master}]
\textbf{GEO-18.} \textbf{Problème.}
\begin{enumerate}
\item[(a)] Démontrez que la surface cubique de Fermat $ x^3 + y^3 + z^3 + w^3 = 0 $ dans
$ \mathbb{P}^3(\mathbb{C}) $ contient exactement 27 droites, et dressez-en la liste.
\item[(b)] Lisez une démonstration du théorème de Cayley--Salmon : \textit{toute} surface cubique lisse sur
$ \mathbb{C} $ contient exactement 27 droites.
\item[(c)] Sur la cubique de Fermat, déterminez combien des 27 droites en rencontrent une donnée.
\end{enumerate}

\end{problem}

\noindent Arthur Cayley montra en 1849 qu'une surface cubique lisse ne porte qu'un nombre
fini de droites, et George Salmon les compta aussitôt : 27, toujours 27 --- le premier grand théorème
énumératif de la géométrie algébrique. La configuration des 27 droites a une symétrie
extraordinaire (son groupe de symétrie est le groupe de Weyl de $ E_6 $), et le résultat a semé
la théorie de l'intersection et le calcul de Schubert. La surface de Fermat rend toute l'histoire
calculable à la main --- rare occasion de toucher du doigt un théorème célèbre par simple
arithmétique.

\begin{refs}
\item Contexte : Surface cubique --- Wikipédia (\url{https://fr.wikipedia.org/wiki/Surface_cubique})
\item Lecture : M. Reid, \textit{Undergraduate Algebraic Geometry}, ch. 7
\end{refs}

\bigskip

\begin{problem}[{Une courbure qu'aucune carte ne peut cacher --- Master, short}]
\textbf{GEO-27.} \textbf{Problème.} Utilisez le Theorema Egregium pour démontrer qu'aucune région d'une sphère, si petite soit-elle, ne peut être envoyée sur le plan sans déformer les distances. Dites exactement quelle quantité fait obstruction.
\end{problem}

\noindent Le \og{} théorème remarquable \fg{} de Gauss affirme que la courbure de Gauss est intrinsèque --- mesurable par un habitant de la surface qui n'en sort jamais --- de sorte qu'une sphère de courbure $ 1/R^2 $ et un plan de courbure nulle ne peuvent jamais être localement isométriques. Toutes les cartes du monde que vous avez vues mentent sur quelque chose, et voici le théorème qui dit qu'elles le doivent.

\begin{refs}
\item Contexte : Theorema egregium --- Wikipédia (\url{https://fr.wikipedia.org/wiki/Theorema_egregium})
\item Lecture : Manfredo do Carmo, \textit{Differential Geometry of Curves and Surfaces}, ch. 4
\end{refs}

\bigskip

\begin{problem}[{Gauss--Bonnet sur un polyèdre --- Master, short}]
\textbf{GEO-28.} \textbf{Problème.} Pour un polyèdre convexe, définissez le défaut angulaire en un sommet comme $ 2\pi $ moins la somme des angles des faces qui s'y rencontrent. Démontrez le théorème de Descartes selon lequel le défaut total sur tous les sommets vaut $ 4\pi $, et expliquez en quoi c'est le théorème de Gauss--Bonnet pour une surface dont la courbure est concentrée en des points.
\end{problem}

\noindent Descartes découvrit cela vers 1630, deux siècles avant Gauss--Bonnet, et c'est l'ombre discrète de ce théorème : la courbure a fui les faces planes et s'est amassée dans les coins, mais le total reste $ 2\pi \chi $. Voir l'ancêtre discret d'un théorème lisse est l'une des expériences les plus satisfaisantes de la géométrie.

\begin{refs}
\item Contexte : Théorème de Descartes sur le défaut angulaire total --- Wikipédia (en anglais) (\url{https://en.wikipedia.org/wiki/Descartes%27_theorem_on_total_angular_defect})
\item Contexte : Formule de Gauss-Bonnet --- Wikipédia (\url{https://fr.wikipedia.org/wiki/Formule_de_Gauss-Bonnet})
\end{refs}

\bigskip

\begin{problem}[{L'inégalité isopérimétrique, démontrée --- Master}]
\textbf{GEO-36.} \textbf{Problème.} GEO-05 montait vers le théorème isopérimétrique et s'arrêtait là où il fallait
l'existence d'un maximiseur. La démonstration de Hurwitz évite entièrement cette lacune en
transformant l'énoncé en analyse de Fourier. Soit $ C $ une courbe fermée simple du plan, de
classe $ C^1 $, de longueur $ L $ et d'aire enfermée $ A $.
\begin{enumerate}
\item[(a)] Démontrez l'inégalité de Wirtinger : si $ f $ est $ 2\pi $-périodique et $ C^1 $ avec
$ \int_0^{2\pi} f = 0 $, alors
\[ \int_0^{2\pi} f(t)^2 \, dt \ \le \ \int_0^{2\pi} f'(t)^2 \, dt , \]
avec égalité exactement lorsque $ f(t) = a\cos t + b \sin t $. (Développez $ f $ en série de
Fourier et appliquez l'identité de Parseval.)
\item[(b)] Démontrez la formule de l'aire $ A = \tfrac12 \oint_C (x \, dy - y \, dx) $ à partir du
théorème de Green, et expliquez pourquoi la convention de signe importe.
\item[(c)] Paramétrez $ C $ sur $ [0, 2\pi] $ proportionnellement à l'abscisse curviligne et
translatez-la de sorte que $ \int_0^{2\pi} x \, dt = 0 $. Démontrez l'inégalité isopérimétrique
\[ L^2 - 4\pi A \ \ge \ 0 , \]
et démontrez que l'égalité force $ C $ à être un cercle.
\item[(d)] Dites précisément où la régularité a servi, et décrivez comment étendre le résultat à toutes
les courbes fermées simples rectifiables par approximation. Comparez ensuite avec l'argument de
symétrisation de Steiner (GEO-05) : qu'apporte chaque démonstration que l'autre n'apporte
pas ?
\end{enumerate}

\end{problem}

\noindent Les arguments de symétrisation de Steiner (1838) étaient beaux et incomplets --- ils
montrent que toute forme autre que le cercle peut être améliorée, mais non qu'une forme optimale
existe, objection que Weierstrass souleva puis leva en développant la méthode directe du calcul des
variations. La démonstration d'Adolf Hurwitz par séries de Fourier, en 1901, contourne entièrement
la difficulté : aucun maximiseur n'est supposé, l'inégalité tombe de l'identité de Parseval, et le
cas d'égalité identifie le cercle en lisant quels coefficients de Fourier survivent. C'est
l'illustration moderne standard d'un principe général --- une identité analytique peut trancher net
un problème géométrique d'extremum --- et sa descendance va de l'inégalité de Brunn--Minkowski aux
théorèmes isopérimétriques sur les variétés.

\begin{refs}
\item Contexte : Théorème isopérimétrique --- Wikipédia (\url{https://fr.wikipedia.org/wiki/Th%C3%A9or%C3%A8me_isop%C3%A9rim%C3%A9trique})
\item Contexte : Inégalité de Wirtinger --- Wikipédia (\url{https://fr.wikipedia.org/wiki/In%C3%A9galit%C3%A9_de_Wirtinger})
\item Lecture : M. do Carmo, \textit{Differential Geometry of Curves and Surfaces}, ch. 1
\item Cours : MIT OCW 18.950 Differential Geometry (\url{https://ocw.mit.edu/courses/18-950-differential-geometry-fall-2008/})
\end{refs}

\bigskip

\begin{problem}[{Courbure totale, et ce que coûte un nœud --- Master}]
\textbf{GEO-37.} \textbf{Problème.} Soit $ \gamma $ une courbe fermée de classe $ C^2 $ dans
$ \mathbb{R}^3 $, paramétrée par l'abscisse curviligne, de tangente unitaire $ T $ et de
courbure $ \kappa $. Son \textit{indicatrice tangentielle} est la courbe sphérique fermée
$ T : [0, L] \to S^2 $, et sa courbure totale est $ \int_\gamma \kappa \, ds $.
\begin{enumerate}
\item[(a)] Démontrez que la courbure totale de $ \gamma $ est égale à la longueur de son indicatrice
tangentielle sur la sphère unité.
\item[(b)] Démontrez le lemme clé : une courbe fermée de $ S^2 $ de longueur inférieure à $ 2\pi $ est
contenue dans un hémisphère ouvert ; et démontrez que l'indicatrice tangentielle d'une courbe
fermée de l'espace ne peut être contenue dans aucun hémisphère ouvert. (Pour la seconde moitié,
intégrez $ T $ le long de la courbe.)
\item[(c)] Déduisez-en le théorème de Fenchel : toute courbe fermée de l'espace a une courbure totale au
moins égale à $ 2\pi $, avec égalité exactement pour les courbes planes convexes.
\item[(d)] Démontrez ou reconstituez soigneusement le théorème de Fáry--Milnor : si $ \gamma $ est
nouée, sa courbure totale dépasse $ 4\pi $. (La voie de Milnor : montrer que la courbure totale
vaut $ 2\pi $ fois le nombre moyen de maxima locaux des fonctions hauteur $ \gamma \cdot v $
sur les vecteurs unitaires $ v $, et qu'une courbe admettant une fonction hauteur à un seul
maximum et un seul minimum est dénouée.)
\end{enumerate}

\end{problem}

\noindent Werner Fenchel démontra la borne $ 2\pi $ en 1929, et Karol Borsuk, qui
l'étendit aux dimensions supérieures en 1947, demanda ce que devait coûter le nouage. István Fáry
répondit en 1949 et John Milnor indépendamment en 1950, alors qu'il était encore étudiant de
premier cycle à Princeton --- l'article qui fit son nom. Le résultat est un pont rare entre les deux
moitiés des pages de géométrie et de topologie de ce site : une quantité de géométrie
différentielle, définie par une intégrale de courbure locale, impose une limite dure à une
propriété purement topologique de la courbe. La borne est optimale au sens où aucune constante plus
petite ne convient, mais elle est loin d'être une égalité --- mesurer la véritable courbure totale
minimale d'un type de nœud donné reste une question vivante.

\begin{refs}
\item Contexte : Théorème de Fary-Milnor --- Wikipédia (\url{https://fr.wikipedia.org/wiki/Th%C3%A9or%C3%A8me_de_Fary-Milnor}), Théorème de Fenchel --- Wikipédia (\url{https://fr.wikipedia.org/wiki/Th%C3%A9or%C3%A8me_de_Fenchel})
\item Article : J. Milnor, \og{} On the total curvature of knots \fg{} (1950), \textit{Annals of Mathematics} 52, 248--257
\item Lecture : M. do Carmo, \textit{Differential Geometry of Curves and Surfaces}, ch. 1 et 5
\item Lecture : C. C. Adams, \textit{The Knot Book}
\end{refs}

\bigskip

\begin{problem}[{Des surfaces faites uniquement de droites --- Master, short}]
\textbf{GEO-38.} \textbf{Problème.} Dites qu'une surface lisse de $ \mathbb{R}^3 $ est
\textit{doublement réglée} si par chacun de ses points passent deux droites distinctes entièrement
contenues dans la surface. Démontrez qu'une surface doublement réglée est contenue dans un plan,
dans un paraboloïde hyperbolique $ z = xy $, ou dans un hyperboloïde à une nappe
$ \frac{x^2}{a^2} + \frac{y^2}{b^2} - \frac{z^2}{c^2} = 1 $.
\end{problem}

\noindent Christopher Wren montra en 1669 que l'hyperboloïde à une nappe est engendré par
une droite mobile, fait qui ressemble à un tour de prestidigitation la première fois qu'on en voit
le modèle. La classification est un classique du XIXe siècle en géométrie analytique des
quadriques, et c'est la raison pour laquelle les tours de refroidissement et les tours en treillis
de Vladimir Choukhov (1896) sont des hyperboloïdes : on obtient une surface à double courbure avec
des barres d'acier rectilignes, bon marché à fabriquer et rigides dans deux directions à la fois.
Notez que les trois surfaces ici ont une courbure de Gauss négative ou nulle --- une droite contenue
dans une surface force la courbure normale à s'annuler dans cette direction, conséquence immédiate
de la machinerie de GEO-13.

\begin{refs}
\item Contexte : Surface réglée --- Wikipédia (\url{https://fr.wikipedia.org/wiki/Surface_r%C3%A9gl%C3%A9e})
\item Contexte : Hyperboloïde --- Wikipédia (\url{https://fr.wikipedia.org/wiki/Hyperbolo%C3%AFde})
\item Lecture : M. do Carmo, \textit{Differential Geometry of Curves and Surfaces}, ch. 3
\end{refs}

\bigskip


\section*{Recherche}

\begin{problem}[{Empilement de sphères : résolu en dimensions 8 et 24, ouvert entre les deux --- Recherche}]
\textbf{GEO-19.} \textbf{Problème.} Avec quelle densité peut-on empiler des boules égales sans recouvrement dans
$ \mathbb{R}^n $ ?
\begin{enumerate}
\item[(a)] Échauffement, à la main : démontrez que l'empilement hexagonal de cercles a pour densité
$ \pi/\sqrt{12} \approx 0{,}9069 $, et démontrez qu'il est optimal parmi les empilements
\textit{en réseau} du plan.
\item[(b)] Problème de compréhension : étudiez la borne de programmation linéaire de Cohn--Elkies et
expliquez comment une unique fonction bien choisie certifie d'un coup une borne supérieure sur la
densité de tout empilement ; lisez ensuite comment Viazovska a construit les fonctions magiques en
dimensions 8 et 24.
\end{enumerate}

\end{problem}

\noindent État des lieux, honnêtement : la dimension 2 est résolue (Thue 1910, complétée
par Fejes Tóth en 1942) ; la dimension 3 est la conjecture de Kepler (1611), démontrée par Hales en
1998 par une analyse de cas massivement informatique et formellement vérifiée par machine en 2017 ;
la dimension 8 a été résolue par Maryna Viazovska en 2016 (le réseau $ E_8 $ est optimal) par une
stupéfiante démonstration de 23 pages via les formes modulaires, étendue en une semaine à la
dimension 24 (le réseau de Leech) par Cohn, Kumar, Miller, Radchenko et Viazovska. Elle a reçu la
médaille Fields en 2022. Toute autre dimension $ n \ge 4 $ reste ouverte --- même $ n = 4 $, et
même si le record est \og{} récemment tombé \fg{} deux fois, personne ne s'attend à ce que les techniques
connues fonctionnent en dimension générale.

\begin{refs}
\item Contexte : Empilement compact (empilement de sphères) --- Wikipédia (\url{https://fr.wikipedia.org/wiki/Empilement_compact})
\item Contexte : Conjecture de Kepler --- Wikipédia (\url{https://fr.wikipedia.org/wiki/Conjecture_de_Kepler})
\item Article : M. Viazovska, \og{} The sphere packing problem in dimension 8 \fg{} (2017), arXiv:1603.04246 (\url{https://arxiv.org/abs/1603.04246})
\item Article : H. Cohn, A. Kumar, S. D. Miller, D. Radchenko, M. Viazovska, \og{} The sphere packing problem in dimension 24 \fg{} (2017), arXiv:1603.06518 (\url{https://arxiv.org/abs/1603.06518})
\end{refs}

\bigskip

\begin{problem}[{Le problème einstein : une seule tuile pour tout désordonner --- Recherche}]
\textbf{GEO-20.} \textbf{Problème.} Un \textit{einstein} (\og{} une pierre \fg{}) est une tuile unique dont les copies pavent
le plan, mais jamais de façon périodique.
\begin{enumerate}
\item[(a)] Échauffement, à la main : démontrez que tout triangle, et en fait tout quadrilatère, convexe ou
non, pave le plan périodiquement --- aucun einstein ne se cache donc parmi eux.
\item[(b)] Problème de compréhension : lisez la démonstration de 2023 selon laquelle le \og{} chapeau \fg{} --- un
polykite à 13 côtés --- est une monotuile apériodique (en utilisant des copies réfléchies), et
rédigez un compte rendu de deux pages de l'architecture de la démonstration : le système de
substitution par métatuiles, et l'argument distinct assisté par ordinateur. Expliquez ce qu'apporte
le \og{} spectre \fg{} : une tuile unique qui est apériodique même sans réflexions.
\end{enumerate}

\end{problem}

\noindent État des lieux, honnêtement : résolu en 2023, et assez récemment pour que cela
ait encore un air d'actualité. Les \textit{ensembles} de tuiles apériodiques remontent à Berger (1966,
plus de 20 000 tuiles de Wang) et à Penrose (1974, deux tuiles) ; savoir si une seule tuile pouvait
suffire est resté ouvert un demi-siècle. En novembre 2022, David Smith, technicien d'imprimerie à
la retraite bricolant des formes, a trouvé le chapeau ; avec Myers, Kaplan et Goodman-Strauss, il
en a démontré l'apériodicité (prépublication en mars 2023), et la même équipe a produit quelques
mois plus tard le \og{} spectre \fg{} chiral. C'est une histoire exemplaire de collaboration entre amateur
et professionnels, et les questions voisines (dimensions supérieures, tuiles à contraintes plus
fortes) demeurent des sujets de recherche actifs.

\begin{refs}
\item Contexte : Problème einstein --- Wikipédia (\url{https://fr.wikipedia.org/wiki/Probl%C3%A8me_einstein})
\item Article : D. Smith, J. S. Myers, C. S. Kaplan, C. Goodman-Strauss, \og{} An aperiodic monotile \fg{} (2023), arXiv:2303.10798 (\url{https://arxiv.org/abs/2303.10798})
\item Article : D. Smith, J. S. Myers, C. S. Kaplan, C. Goodman-Strauss, \og{} A chiral aperiodic monotile \fg{} (2023), arXiv:2305.17743 (\url{https://arxiv.org/abs/2305.17743})
\end{refs}

\bigskip

\begin{problem}[{Le problème du déménagement de canapé --- Recherche}]
\textbf{GEO-21.} \textbf{Problème.} Quelle est la plus grande aire d'une forme plane que l'on puisse déplacer (par
translations et rotations) autour d'un angle droit dans un couloir de largeur unité ?
\begin{enumerate}
\item[(a)] À la main : vérifiez qu'un carré unité passe, qu'un demi-disque de rayon 1 (d'aire
$ \pi/2 $) passe, et que la forme de Hammersley --- un rectangle amputé d'une échancrure
semi-circulaire, flanqué de deux quarts de disque, d'aire $ \pi/2 + 2/\pi \approx 2{,}2074 $ ---
tourne le coin.
\item[(b)] Démontrez une borne supérieure : montrez que l'aire d'un canapé est au plus
$ 2\sqrt{2} $, ou mieux.
\item[(c)] Problème de compréhension : décrivez le canapé en 18 morceaux de Gerver et ce que
signifierait son optimalité.
\end{enumerate}

\end{problem}

\noindent État des lieux, honnêtement : posé par Leo Moser en 1966 ; Gerver construisit en
1992 son canapé à 18 arcs, d'aire $ \approx 2{,}2195 $, et le conjectura optimal. En novembre
2024, Jineon Baek (université Yonsei) a mis en ligne une démonstration de 119 pages, \og{} Optimality of
Gerver's Sofa \fg{}, prétendant résoudre le problème --- notamment sans assistance informatique. En 2026,
la démonstration est toujours en cours de relecture par les pairs, les experts se montrant
optimistes quant à sa correction (Scientific American l'a comptée parmi les percées de 2025),
mais il faut honnêtement la dire \og{} annoncée, en cours d'évaluation \fg{} plutôt que réglée. La variante
ambidextre --- un canapé qui tourne à la fois les coins à gauche et à droite, avec l'optimum conjecturé
par Romik --- reste ouverte.

\begin{refs}
\item Contexte : Problème du sofa --- Wikipédia (\url{https://fr.wikipedia.org/wiki/Probl%C3%A8me_du_sofa})
\item Article : J. Baek, \og{} Optimality of Gerver's Sofa \fg{} (2024), arXiv:2411.19826 (\url{https://arxiv.org/abs/2411.19826})
\end{refs}

\bigskip

\begin{problem}[{Le problème de Plateau et ses affaires inachevées --- Recherche}]
\textbf{GEO-22.} \textbf{Problème.} Le problème de Plateau demande : toute courbe de bord fermée de l'espace
sous-tend-elle une surface d'aire minimale ?
\begin{enumerate}
\item[(a)] Problème de compréhension : expliquez pourquoi la solution de Douglas--Radó des années 1930
(surfaces \textit{de type disque} d'aire minimale) ne peut rendre compte des films de savon réels,
comme le film sur un cadre de fil tétraédrique, qui présente des courbes singulières où trois
nappes se rencontrent à $ 120^\circ $ et des points où quatre de ces courbes se rencontrent à
$ \approx 109{,}47^\circ $ --- les lois de Plateau, observées expérimentalement dans les années
1870 et démontrées par Jean Taylor en 1976.
\item[(b)] Formulez précisément, et comparez, deux formulations modernes non équivalentes de \og{} surface
sous-tendue par une courbe \fg{} (par exemple les courants de Federer--Fleming face à une formulation
ensembliste, de type film de savon), et expliquez ce que \og{} solution \fg{} veut dire dans chacune.
\end{enumerate}

\end{problem}

\noindent État des lieux, honnêtement : le problème de type disque a été résolu
indépendamment par Jesse Douglas et Tibor Radó vers 1930--31, et Douglas en reçut la première
médaille Fields (1936). Mais choisir les bonnes définitions s'est révélé le problème le plus
profond, donnant naissance à la théorie géométrique de la mesure. Beaucoup reste véritablement
ouvert : la structure des singularités des surfaces minimisant l'aire en dimension et codimension
élevées (le monumental théorème de régularité partielle d'Almgren borne l'ensemble singulier, mais
sa structure fine est une frontière active), et les questions d'existence de type Plateau pour les
modèles de films de savon continuent d'alimenter la recherche aujourd'hui.

\begin{refs}
\item Contexte : Problème de Plateau --- Wikipédia (\url{https://fr.wikipedia.org/wiki/Probl%C3%A8me_de_Plateau})
\item Contexte : Conditions de Plateau (lois de Plateau) --- Wikipédia (\url{https://fr.wikipedia.org/wiki/Conditions_de_Plateau})
\item Contexte : Théorie géométrique de la mesure --- Wikipédia (\url{https://fr.wikipedia.org/wiki/Th%C3%A9orie_g%C3%A9om%C3%A9trique_de_la_mesure})
\item Lecture : M. do Carmo, \textit{Differential Geometry of Curves and Surfaces}, ch. 3
\end{refs}

\bigskip

\begin{problem}[{Que sait-on de l'empilement de sphères ? --- Recherche, short}]
\textbf{GEO-29.} \textbf{Problème.} Dressez un tableau précis de l'empilement de sphères le plus dense connu en dimensions $ 1 $ à $ 24 $, en indiquant pour chaque dimension si l'optimalité est \textit{démontrée} ou seulement \textit{conjecturée}. Seules les dimensions $ 1, 2, 3, 8 $ et $ 24 $ sont réglées.
\end{problem}

\noindent Composer ce tableau vous oblige à mesurer combien nos connaissances sont minces : trois des cinq cas résolus ont exigé soit une démonstration par ordinateur, soit les formes modulaires, et la dimension $ 4 $ --- que vous pouvez presque vous représenter --- est ouverte. C'est aussi un utile antidote à l'idée que les problèmes en petite dimension seraient forcément faciles.

\begin{refs}
\item Contexte : Empilement compact (empilement de sphères) --- Wikipédia (\url{https://fr.wikipedia.org/wiki/Empilement_compact})
\item Voir aussi : Problèmes ouverts (\url{open-problems.html})
\end{refs}

\bigskip

\begin{problem}[{L'aiguille de Kakeya, et ce qu'elle coûte --- Recherche, short}]
\textbf{GEO-30.} \textbf{Problème.} Expliquez la construction de Besicovitch (1919) montrant qu'une aiguille de longueur unité peut être tournée de $ 180^\circ $ à l'intérieur d'un ensemble d'aire arbitrairement petite, puis énoncez précisément ce qu'affirme au contraire la conjecture de Kakeya --- qu'un tel ensemble doit néanmoins avoir une dimension de Hausdorff maximale. Notez l'état des lieux : démontrée en dimension trois par Wang et Zahl en 2025, ouverte en dimension quatre et au-delà.
\end{problem}

\noindent La tension est ici tout l'enjeu : la mesure peut s'évanouir alors que la dimension ne le peut pas, et apprendre à tenir séparées ces deux notions de taille est un vrai pas conceptuel. Ce problème est aussi directement lié à la conjecture de restriction en analyse harmonique, ce qui explique pourquoi les analystes s'intéressent à une question qui sonne comme de la géométrie récréative.

\begin{refs}
\item Contexte : Ensemble de Besicovitch (ensemble de Kakeya) --- Wikipédia (\url{https://fr.wikipedia.org/wiki/Ensemble_de_Besicovitch})
\item Voir aussi : Problèmes ouverts (\url{open-problems.html}), Analyse réelle (\url{pure-real-analysis.html})
\end{refs}

\bigskip

\begin{problem}[{Combien de boules peuvent en toucher une seule ? --- Recherche, short}]
\textbf{GEO-39.} \textbf{Problème.} (Ouvert dans presque toutes les dimensions.) Le nombre de
contacts $ \tau_n $ est le plus grand nombre de boules unités sans recouvrement de
$ \mathbb{R}^n $ pouvant toutes toucher une boule unité fixée. Déterminez $ \tau_n $ --- les
valeurs exactes ne sont connues que pour $ n = 1, 2, 3, 4, 8, 24 $, et le plus petit cas ouvert
est $ n = 5 $, où tout ce que l'on sait est $ 40 \le \tau_5 \le 44 $.
\end{problem}

\noindent Newton et David Gregory se disputèrent en 1694 pour savoir si douze ou treize
boules pouvaient en toucher une treizième dans l'espace ordinaire ; Newton disait douze et avait
raison, mais les jeux dans la disposition à douze boules sont assez larges pour qu'aucune
démonstration ne paraisse avant Schütte et van der Waerden en 1953. Musin régla $ n = 4 $ en 2003
en poussant plus loin la méthode de programmation linéaire de Delsarte, et $ \tau_8 = 240 $ et
$ \tau_{24} = 196560 $ tombèrent indépendamment sous Odlyzko--Sloane et Levenshtein en 1979, avec
ces mêmes réseaux $ E_8 $ et de Leech qui résoudront plus tard le problème d'empilement de
GEO-19. Notez que ce n'est pas ce problème-là : le contact est une question locale sur les voisins
d'une boule, la densité une question globale, et aucune des deux réponses n'entraîne l'autre.

\begin{refs}
\item Contexte : Nombre de contacts --- Wikipédia (\url{https://fr.wikipedia.org/wiki/Nombre_de_contacts})
\item Contexte : Réseau E8 --- Wikipédia (en anglais) (\url{https://en.wikipedia.org/wiki/E8_lattice}), Réseau de Leech --- Wikipédia (\url{https://fr.wikipedia.org/wiki/R%C3%A9seau_de_Leech})
\item Lecture : J. H. Conway et N. J. A. Sloane, \textit{Sphere Packings, Lattices and Groups}, ch. 1
\end{refs}

\bigskip

\begin{problem}[{Un carré sur toute courbe fermée ? --- Recherche}]
\textbf{GEO-40.} \textbf{Problème.} (Ouvert.) Une \textit{courbe de Jordan} est l'image d'une application continue
injective $ S^1 \to \mathbb{R}^2 $. Toeplitz demanda en 1911 si toute courbe de Jordan contient
les quatre sommets d'un carré.
\begin{enumerate}
\item[(a)] Démontrez la version rectangle, qui n'exige aucune régularité. L'espace des paires non
ordonnées de points de $ S^1 $ est un ruban de Möbius dont le bord est la diagonale ; envoyez une
paire non ordonnée $ \{p, q\} $ de points de la courbe sur
$ (\,m(p,q), \ |p - q|\,) \in \mathbb{R}^2 \times [0,\infty) $, où $ m $ est le milieu, et
montrez qu'un défaut d'injectivité est exactement un rectangle inscrit dans la courbe --- puis
montrez que l'injectivité est impossible, car le ruban de Möbius ne se plonge pas dans un
demi-espace avec son bord dans le plan frontière.
\item[(b)] Démontrez la version carré pour les courbes convexes, et lisez la démonstration d'Emch (1916)
pour les courbes analytiques par morceaux.
\item[(c)] Étudiez le théorème de Greene et Lobb (2020) : toute courbe de Jordan lisse inscrit un
rectangle de \textit{tout} rapport d'aspect. Expliquez comment le problème est reformulé en une
question sur deux tores lagrangiens de $ \mathbb{C}^2 $, et ce que la géométrie symplectique
apporte que la topologie plane ne pouvait pas apporter.
\item[(d)] Le problème lui-même : démontrez que toute courbe de Jordan inscrit un carré, ou exhibez-en une
qui n'en inscrit pas. Dites précisément quelle étape des démonstrations connues casse pour une
courbe seulement continue.
\end{enumerate}

\end{problem}

\noindent État des lieux, honnêtement : ouvert. Otto Toeplitz posa la question en 1911, et
un siècle de travaux l'a réglée pour toute classe raisonnable de courbes --- courbes convexes,
courbes analytiques par morceaux (Emch, 1916), courbes localement monotones (Stromquist, 1989),
courbes présentant diverses symétries --- tandis que le cas continu général résiste, parce qu'une
courbe de Jordan peut être assez sauvage pour qu'aucun argument de passage à la limite ne survive.
Le progrès récent le plus frappant est le théorème du piquet rectangulaire de Joshua Greene et
Andrew Lobb (\textit{Annals of Mathematics} 194, 2021), qui importe la géométrie symplectique dans un
problème qui semble relever d'un premier cours, et qu'ils ont étendu peu après aux quadrilatères
inscriptibles inscrits. En 2026, on ne connaît ni contre-exemple ni démonstration générale.

\begin{refs}
\item Contexte : Problème du carré inscrit --- Wikipédia (\url{https://fr.wikipedia.org/wiki/Probl%C3%A8me_du_carr%C3%A9_inscrit})
\item Contexte : Théorème de Jordan --- Wikipédia (\url{https://fr.wikipedia.org/wiki/Th%C3%A9or%C3%A8me_de_Jordan})
\item Article : J. E. Greene et A. Lobb, \og{} The rectangular peg problem \fg{} (2021), \textit{Annals of Mathematics} 194, arXiv:2005.09193 (\url{https://arxiv.org/abs/2005.09193})
\item Voir aussi : Problèmes ouverts (\url{open-problems.html}), Topologie (\url{pure-topology.html})
\end{refs}

\bigskip


\section*{Pour aller plus loin}


\vfill
\noindent\emph{Hors de la Caverne} --- des probl\`emes, pas de r\'eponses.
\end{document}
